% =====================================================================
%  附录：审校批注索引
%  由工具从正文各分片中的 % [Tutor Note: ...] 自动汇总生成；
%  源码中的批注为 LaTeX 注释，编译后不显示，此处以正文形式列出。
% =====================================================================
\newpage
\section*{附录　审校批注索引}
\addcontentsline{toc}{section}{附录　审校批注索引}
\markboth{附录　审校批注索引}{}
\noindent\small
本附录把正文源码中的 \verb|% [Tutor Note: ...]| 批注按章节顺序汇总，便于不读源码时检索。
每条批注前标出所属小节；若批注本身注明页码，则在括注中给出原文电子底本的页码
（原文印刷页号 $=n-1$）。
\normalsize

\par\medskip\noindent\hangindent=3.2em\hangafter=1 \textbf{[1]}~\textbf{引言}~见原文 PDF 第 2 页引言第一段。文中“Trieste Summer School”指 1995 年在意大利的里雅斯特（Trieste）举办的暑期学校系列讲座，本讲义即这些讲座的整理稿；原文称其为“very late written version”（很晚才写成的版本），故译作“成文远晚于讲座本身”。
\par
\par\medskip\noindent\hangindent=3.2em\hangafter=1 \textbf{[2]}~\textbf{引言}~见原文 PDF 第 3 页引言第二段。原文称含物质时“the duality group is no longer pure $SL(2,Z)$”，其中“pure”指对偶群不再是整个 $SL(2,Z)$，而是它的真子群（同余子群）；具体形式与物质场的味数 $N_f$ 有关，详见原文第 4.6 节（原文 PDF 第 75 页起）。
\par
\par\medskip\noindent\hangindent=3.2em\hangafter=1 \textbf{[3]}~\textbf{1.1}~见原文 PDF 第 3 页式 (1)。此处采用的是 $c=\hbar=1$ 且库仑定律带因子 $(4\pi)^{-1}$ 的“有理化”单位制（Lorentz–Heaviside 型），故点电荷场强为 $\vec{E}=q\vec{r}/4\pi r^{3}$ 而非 $q\vec{r}/r^{3}$。后文的 $4\pi$、$e=q/\hbar$ 以及磁荷 $g$ 的归一化全都由它确定；若换用 SI 或高斯制，电荷量子化条件会整体差一个 $4\pi$，核对时请务必以此单位制为准。
\par
\par\medskip\noindent\hangindent=3.2em\hangafter=1 \textbf{[4]}~\textbf{1.1}~见原文 PDF 第 3 页式 (2) 与第 4 页式 (3)。在号差 $(+,-,-,-)$、$\varepsilon^{0123}=+1$ 下，$F_{0i}=E^i$ 但升指标得 $F^{0i}=-E^i$，同理 $\widetilde F^{0i}=-B^i$；式 (3) 的两个矩阵里 $E$、$B$ 的符号分布正是这一约定的结果，此处是最容易漏掉负号的地方，请逐项核对。
\par
\par\medskip\noindent\hangindent=3.2em\hangafter=1 \textbf{[5]}~\textbf{1.1}~见原文 PDF 第 4 页式 (5)、(6)。要点有二：其一，对偶不变性只在场与流“同时”按式 (6) 变换时才成立，单独变换场并不是麦克斯韦方程组的对称性；其二，原文所说连续的 $SO(2)$ 群指 $(F,\widetilde F)$ 平面内的转动（等价于对复组合 $F+i\widetilde F$ 作相位变换），这正是后文 $S$ 对偶连续版本的雏形。
\par
\par\medskip\noindent\hangindent=3.2em\hangafter=1 \textbf{[6]}~\textbf{1.2}~见原文 PDF 第 5 页式 (8)。第二个等号是分部积分：$\partial_j\hat{x}^i=(\delta_{ij}-x_ix_j/r^2)/r$，再把 $E^j\partial_j\hat{x}^i$ 写成 $\partial_j(E^j\hat{x}^i)-\hat{x}^i\partial_j E^j$。原文只说“球对称使第一个积分消失”，对真正给出式 (9) 的 $\vec{\nabla}\cdot\vec{E}=q\delta^{3}(\vec{r})$ 一项轻轻带过——恰恰是原点处的 $\delta$ 源贡献了 $-qg/4\pi\,\hat{r}$。这是本段最容易“差一个符号”的地方，请按 $\partial_j\hat{x}^i$ 的符号重新核对。
\par
\par\medskip\noindent\hangindent=3.2em\hangafter=1 \textbf{[7]}~\textbf{1.2}~见原文 PDF 第 6 页式 (10) 前一句。原文“if we assume that orbital angular momentum is quantized. Then it follows that”在句号处应作逗号（原文标点笔误），译文的句式已按正确逻辑重写。物理上这里假设的是“总角动量（含场角动量 $\vec{J}_{em}$）以 $n\hbar/2$ 量子化”，而非纯轨道角动量；随后原文承认这与两个组元都是玻色子相矛盾，并把解释归于文献 [5,6]（Goldhaber、Wilczek）。下一小节由狄拉克弦给出的推导不依赖此假设。
\par
\par\medskip\noindent\hangindent=3.2em\hangafter=1 \textbf{[8]}~\textbf{1.2}~见原文 PDF 第 6 页式 (13)。注意 $n_{12}$ 依赖两粒子的编号顺序，交换 $1\leftrightarrow2$ 时 $n_{12}\to-n_{12}$；而 $q_1g_2-q_2g_1$ 正是二维矢量 $(q_i,g_i)$ 的行列式，在 $(q_i,g_i)$ 整体 $SO(2)$ 转动下不变。这一“量子化条件与连续对偶相容”的事实，正是后文（§4）把对偶群视为 $SL(2,Z)$ 的子群、并研究其单值变换的出发点。
\par
\par\medskip\noindent\hangindent=3.2em\hangafter=1 \textbf{[9]}~\textbf{1.3}~见原文 PDF 第 7 页第 2 行。原文写 $\vec{B}_{mon}=\vec{\nabla} \vec{A}_{mon}$，漏掉叉乘号，应为 $\vec{B}_{mon}=\vec{\nabla}\times\vec{A}_{mon}$（原文笔误，译文照录）。正因为磁单极场满足 $\vec{\nabla}\cdot\vec{B}_{mon}\ne 0$，这样的 $\vec{A}_{mon}$ 才无法整体存在。
\par
\par\medskip\noindent\hangindent=3.2em\hangafter=1 \textbf{[10]}~\textbf{1.3}~见原文 PDF 第 7 页式 (14)。其中 $g/4\pi r^{2}$ 一项应为 $g\hat{r}/4\pi r^{2}$（即 $\vec{B}_{mon}=g\vec{r}/4\pi r^{3}$），原文漏写单位矢量 $\hat{r}$（原文笔误，译文照录）；补上它，$\vec{B}_{mon}$ 的方向信息才完整，也才能与 $\vec{\nabla}\cdot\vec{B}_{mon}=g\delta^{3}(\vec{r})$ 自洽。
\par
\par\medskip\noindent\hangindent=3.2em\hangafter=1 \textbf{[11]}~\textbf{1.3}~见原文 PDF 第 7 页“Aharonov–Bohm 相位”一段，紧接式 (14) 之后。相位 $e\oint_\Gamma\vec{A}_{sol}\cdot d\vec{l}=eg$ 中的 $e$ 是经 $\hbar$ 归一化的耦合常数 $e=q/\hbar$（原文如此，并非裸电荷），故 $eg=2\pi n$ 与上小节的 $qg/4\pi\hbar=n/2$ 完全等价；两种推导的一致性全靠这一归一化，切勿把 $e$ 当作 $q$ 代入。
\par
\par\medskip\noindent\hangindent=3.2em\hangafter=1 \textbf{[12]}~\textbf{1.4}~见原文 PDF 第 8 页“Derrick 定理”一句。原文只给结论，其证明是标度论证：对静态场作标度变换 $\phi_\lambda(\vec{x})=\phi(\lambda\vec{x})$，能量泛函可写成 $E(\lambda)=\lambda^{2-d}K+\lambda^{-d}V$（$K$ 为梯度项、$V$ 为势能项，均非负）；$d>1$ 时 $\mathrm{d}E/\mathrm{d}\lambda=0$ 无解，故 $d\ge2$ 的标量场论不存在静态有限能量解。原文另把这归因于孤立子在远距离处的非平凡结构，这是从位形角度的直观说法。
\par
\par\medskip\noindent\hangindent=3.2em\hangafter=1 \textbf{[13]}~\textbf{1.4}~见原文 PDF 第 8 页式 (15)、(16)。$W_\mu=W^a_\mu T^a$ 是矩阵值场，$G_{\mu\nu}=G^a_{\mu\nu}T^a$；用 $(T^a)_{bc}=-if^a_{bc}$ 与 $f^{abc}=\epsilon^{abc}$ 展开 $ie[W_\mu,W_\nu]$，得 $G^a_{\mu\nu}=\partial_\mu W^a_\nu-\partial_\nu W^a_\mu-e\epsilon^{abc}W^b_\mu W^c_\nu$。式中的因子 $i$ 与 $D_\mu\phi^a$ 里的 $-e$ 是配套的，符号弄反会直接改变式 (18) 的运动方程，并影响后文 BPS 界与磁荷的符号。
\par
\par\medskip\noindent\hangindent=3.2em\hangafter=1 \textbf{[14]}~\textbf{1.4}~见原文 PDF 第 9 页微扰谱表格。表中“Spin”一列对光子与 $W^\pm$ 填 $\hbar$、对 $H$ 填 $0$，这是原文的写法：自旋一栏理应填自旋量子数 $1,1,0$（即 $s\hbar$ 中的 $s$），原文直接填入 $\hbar$，单位上亦不自洽，属原文笔误；为与源公式块保持一致，表头 Mass/Spin/Charge 保留原文，译名即“质量/自旋/电荷”。表中 $W^\pm$ 的质量 $ae\hbar=aq$ 与电荷 $\pm q=\pm e\hbar$ 彼此一致。
\par
\par\medskip\noindent\hangindent=3.2em\hangafter=1 \textbf{[15]}~\textbf{1.4}~本片原文拼写笔误汇总（译文均按正确词义翻译，源文件未改）：p2 “intented”→intended；p3 “proporties”→properties、“Minkowsky”→Minkowski；p4 “quations”→equations；p5 “elecromagnetic”→electromagnetic；p6 “fermios”→fermions，且“both of the components are bosonic.”句首小写；p7 “intruduced”→introduced；p8 “satifying”→satisfying、“representaion”→representation、“cavariant”→covariant。另有句读笔误一处，见式 (10) 前的批注。
\par
\par\medskip\noindent\hangindent=3.2em\hangafter=1 \textbf{[16]}~\textbf{1.5}~见原文 PDF 第 9 页开头。有限能量条件为 $E=\int\theta_{00}\,d^3x<\infty$：三维体积元带来 $r^2dr$，故要求 $\theta_{00}$ 的衰减快于 $1/r^3$（而不是快于 $1/r^2$）；这里的 $\theta_{00}$ 在静态情形下就是能量密度。
\par
\par\medskip\noindent\hangindent=3.2em\hangafter=1 \textbf{[17]}~\textbf{1.5}~见原文 PDF 第 10 页式 (24)。此处系数最容易写错：第一项的 $1/(ea^2)$ 中 $a^2$ 来自 $\vec{\phi}_{vac}\cdot\vec{\phi}_{vac}=a^2$，第二项的 $1/a$ 保证 $A_\mu$ 是归一化的 $U(1)$ 规范场。代回式 (23) 时只需两个恒等式：$\vec{\phi}_{vac}\cdot\partial_\mu\vec{\phi}_{vac}=0$ 与 $(\vec{\phi}_{vac}\times\partial_\mu\vec{\phi}_{vac})\times\vec{\phi}_{vac}=a^2\partial_\mu\vec{\phi}_{vac}$。
\par
\par\medskip\noindent\hangindent=3.2em\hangafter=1 \textbf{[18]}~\textbf{1.5}~原文笔误，见原文 PDF 第 10 页式 (26) 之后。“Since $\phi^2=a$;”有两处问题：(1) 应为 $\phi^2=a^2$，与式 (23) 一致；(2) 分号应为逗号。译文照录公式原样，只按文意断句。
\par
\par\medskip\noindent\hangindent=3.2em\hangafter=1 \textbf{[19]}~\textbf{1.5}~见原文 PDF 第 10 页式 (28)。原文说它与狄拉克条件“相差因子 2”，核对如下：取 $q_0=e\hbar/2$，则 $g_\Sigma q_0/(4\pi\hbar)=-\nu/2$，而狄拉克条件为 $qg/(4\pi\hbar)=n/2$，故要求 $\nu=-n$。式 (26) 中的负号来自 $B^i=-\frac{1}{2}\epsilon^{ijk}F_{jk}$ 与曲面定向的约定，不要与式 (27) 中 $\nu$ 的符号约定混为一谈。
\par
\par\medskip\noindent\hangindent=3.2em\hangafter=1 \textbf{[20]}~\textbf{1.5}~见原文 PDF 第 11 页式 (29) 之后。核对边界条件：$r\to0$ 时 $H\sim\xi$ 使 $\phi^a=(x^a/er^2)H\to c_1x^a\to0$，故场在原点正则；$r\to\infty$ 时 $H\sim\xi$ 给出 $\phi^a\to ax^a/r=a\hat{x}^a$，而 $K\to0$ 给出 $W^a_i\to-\epsilon^a_{ij}x^j/er^2$。注意 $H$、$K$ 的自变量是 $\xi=aer$ 而不是 $r$，做量纲核对时不要丢掉 $e$ 与 $a$。
\par
\par\medskip\noindent\hangindent=3.2em\hangafter=1 \textbf{[21]}~\textbf{1.5}~原文 PDF 第 11 页这一段有数处拼写笔误：“easiy”应为“easily”，“guage”应为“gauge”，“treate”应为“treat”，“depends of”应为“depends on”；译文均按原意处理。关于集体坐标论证：静态双子解带有一个连续参数（即 $W^a_0$ 的系数），它对应于保持希格斯真空的 $U(1)$ 规范方向的零模；把该规范自由度作为集体坐标量子化后，其共轭动量（电荷）取分立值 $q=n\hbar e$，细节见文献 [12,13]。
\par
\par\medskip\noindent\hangindent=3.2em\hangafter=1 \textbf{[22]}~\textbf{1.6}~见原文 PDF 第 12 页式 (30)。把面积分化为体积分用了两个不同方程：磁荷用非阿贝尔比安基恒等式 $D_i{\cal B}^a_i=0$，电荷用高斯定律（即 (\ref{XXIX}) 的第一个方程）。注意等号链中“面积分”与“体积分”之间的那一步并不只是把 $a$ 提出：需要把 $\phi^a/a$ 作为单位矢量场、分部积分并用 $D_i{\cal B}^a_i=0$，才能把被积函数中的 $\phi^a$ 换成 $(D^i\phi)^a$；这是本节最容易被读者一带而过的地方。
\par
\par\medskip\noindent\hangindent=3.2em\hangafter=1 \textbf{[23]}~\textbf{1.6}~易错预警，见原文 PDF 第 12 页式 (31)。平方配方时务必与式 (30) 中 $g,q$ 的定义对照：把方括号内展开，除保留 $\frac12[({\cal E}^a_k)^2+({\cal B}^a_k)^2+(D_k\phi^a)^2+(D_0\phi^a)^2]$ 外，交叉项为 $-\sin\theta\,{\cal E}\cdot D\phi-\cos\theta\,{\cal B}\cdot D\phi=-a(q\sin\theta+g\cos\theta)$（注意是负号！），移项后才得到式 (31) 末行的 $+a(q\sin\theta+g\cos\theta)$。若把交叉项符号弄反，就会得出 $M\le a\sqrt{q^2+g^2}$ 的错误结论。对 $\theta$ 求极值时用 $\tan\theta=q/g$，最大值为 $\sqrt{q^2+g^2}$。
\par
\par\medskip\noindent\hangindent=3.2em\hangafter=1 \textbf{[24]}~\textbf{1.6}~见原文 PDF 第 12 页。此处对 $\hbar$ 的处理较随意：由 $M_W=ae\hbar=aq$ 与 $\nu/\alpha$ 的写法可反推出 $\alpha\equiv q^2/(4\pi\hbar)$，只有取 $\hbar=1$ 时才与通常的 $\alpha=e^2/4\pi\sim1/137$ 相符。原文此处只作定性结论（磁单极远重于 W 玻色子），不必纠结 $\hbar$ 的幂次；$\nu=1$ 或 $1/4$ 则取决于把“电子电荷”取作 $q$ 还是 $q/2$。
\par
\par\medskip\noindent\hangindent=3.2em\hangafter=1 \textbf{[25]}~\textbf{1.6}~见原文 PDF 第 13 页。式 (32) 其实是两类一阶（BPS）方程的合写：$D_0\phi=0$ 与 $\vec{\cal E}=\sin\theta\,D\phi$、$\vec{\cal B}=\cos\theta\,D\phi$；把 't Hooft–Polyakov 拟设代入这些一阶方程即可精确求解（原文引 [14,15]）。另注意：$\lambda\to0$ 时 $V=0$，$a$ 不再由势的极小值固定，故 $|\phi|=a$ 只能作为边界条件强加，这一点在后文 Seiberg–Witten 分析中至关重要。
\par
\par\medskip\noindent\hangindent=3.2em\hangafter=1 \textbf{[26]}~\textbf{1.7}~术语与记号说明，见原文 PDF 第 13 页。原文用同一个字母 $g$ 表示两件完全不同的东西：群元素 $g$（如式 (36) 的 $g(s,t)$）与磁荷 $g$（见式 (30) 与式 (35)）。阅读时须按上下文区分，切勿把 $g(s,t)$ 的指标结构套到磁荷上。
\par
\par\medskip\noindent\hangindent=3.2em\hangafter=1 \textbf{[27]}~\textbf{1.7}~易错预警，见原文 PDF 第 14 页式 (35)。与 §1.6 的 $M\ge a\sqrt{g^2+q^2}$ 相比这里多出因子 $\sqrt2$：式 (33) 的希格斯动能是复标量的 $(D_\mu\phi)^\dagger D^\mu\phi$，平方配方时须写成 $\frac{1}{2}|{\cal E}-\sqrt2\,D\phi\sin\theta|^2$、$\frac{1}{2}|{\cal B}-\sqrt2\,D\phi\cos\theta|^2$ 之类的形式（实部、虚部各贡献一份），交叉项才给出 $\sqrt2\,a(q\sin\theta+g\cos\theta)$；若照搬 §1.6 实伴随场的 $\frac12(D\phi)^2$ 配方，就会漏掉这个 $\sqrt2$。
\par
\par\medskip\noindent\hangindent=3.2em\hangafter=1 \textbf{[28]}~\textbf{1.7}~含糊点拨，见原文 PDF 第 14 页。原文“$G_{\mu\nu}$ 取值于 $G$ 的某个子群”严格说应为“取值于该子群的李代数”；这里只是省略了“李代数”三字，物理内容不变。又：原文此处有重复词 “the the little group”，系笔误。
\par
\par\medskip\noindent\hangindent=3.2em\hangafter=1 \textbf{[29]}~\textbf{1.7}~原文笔误，见原文 PDF 第 14 页：“from this square to the sphare” 中 “sphare” 应为 “sphere”。
\par
\par\medskip\noindent\hangindent=3.2em\hangafter=1 \textbf{[30]}~\textbf{1.7}~图 2 的读法，见原文 PDF 第 15 页。$(s,t)$ 正方形的整个边界都映到基点 $r_0$，故对每个固定的 $s$，$t\in[0,1]$ 是 $\Sigma$ 上一条由 $r_0$ 出发又回到 $r_0$ 的闭曲线；右图中这些闭曲线按同伦类分层，其绕数即磁荷。注意左图中竖直箭头只表示 $t$ 增大方向，图中的“曲线族”是同一映射 $r(s,t)$ 在不同 $s$ 处的像。
\par
\par\medskip\noindent\hangindent=3.2em\hangafter=1 \textbf{[31]}~\textbf{1.7}~见原文 PDF 第 15 页式 (36)。$P$ 表示路径编序（沿 $t$ 增大方向）；指数中 $ie\int\vec{W}\cdot d\vec{r}$ 的量纲须与式 (33) 的 $D_\mu=\partial_\mu+ieW_\mu$ 一致，核对时勿额外添加 $e$ 或 $2\pi$。
\par
\par\medskip\noindent\hangindent=3.2em\hangafter=1 \textbf{[32]}~\textbf{1.7}~含糊点拨，见原文 PDF 第 16 页的正合序列。原文“$\pi_2(G)$ 可作为单位元嵌入 $\pi_2(G/H)$”的措辞不严谨：对紧致连通 $G$ 有 $\pi_2(G)=0$，于是由正合性得 $\pi_2(G/H)$ 单射地映入 $\pi_1(H)$，其像正是（在 $G$ 中同伦于平庸路径的）$\pi_1(H)_G$，故 $\pi_2(G/H)\cong\pi_1(H)_G$。原文所谓“嵌入”只是说明 $\pi_2(G)=0$ 这一环在序列中不贡献，读者不必按字面理解为群嵌入。完整讨论见文献 [2]（Coleman）。
\par
\par\medskip\noindent\hangindent=3.2em\hangafter=1 \textbf{[33]}~\textbf{1.7}~见原文 PDF 第 16 页。原文称“后面我们会证明 $g=4\pi N/e$”，但本分片后文并未再给出独立证明；读者可对照式 (37)–(42) 的广义量子化条件自行补出（比较 $h(t)$ 的指数中 $4\pi N/a$ 的归一化与式 (41) 中 $G_0$ 前的 $ie/4\pi$）。另注意半整数 $N$ 的路径在 $SO(3)$ 中虽闭合，却生成 $\pi_1(SO(3))={\bf Z}_2$ 的非平凡类，不在 $\pi_1(H)_G$ 中，故最终两种情形都只保留整数 $N$。
\par
\par\medskip\noindent\hangindent=3.2em\hangafter=1 \textbf{[34]}~\textbf{1.8}~含糊点拨，见原文 PDF 第 17 页。原文“the magnetic flux in a solid angle subtended at the origin by the path”措辞较绕：应理解成“穿过该路径所张立体角（以原点为顶点）的磁通量”，即 $\Omega(s)$ 本身是立体角、乘以 $B$ 才是通量；这正是 $\alpha(s)=d\Omega/ds$ 的几何含义。
\par
\par\medskip\noindent\hangindent=3.2em\hangafter=1 \textbf{[35]}~\textbf{1.8}~含糊点拨，见原文 PDF 第 18 页。原文“If all values of $n$ are allowed, then $U(1)\cap K\subset Z_N$”一步较绕：其含义是整体上 $H$ 不是直积 $U(1)\times K$，而是 $(U(1)\times K)/Z_N$ 型的结构，故 $U(1)$ 与 $K$ 的中心在整体上重合；对色单态（$k=1$）回到通常的狄拉克条件，而对色非单态则得到以 $1/N$ 为单位的分数电荷。注意原文用 $|s>$、$|c>$ 而不是 $|s\rangle$、$|c\rangle$，系作者笔误，译文照录。
\par
\par\medskip\noindent\hangindent=3.2em\hangafter=1 \textbf{[36]}~\textbf{1.9}~原文笔误，见原文 PDF 第 18 页：“the magnetic componenets”中 “componenets” 应为 “components”。
\par
\par\medskip\noindent\hangindent=3.2em\hangafter=1 \textbf{[37]}~\textbf{1.9}~见原文 PDF 第 18–19 页。由式 (42) 得到 $e\beta^a\omega^a\in2\pi{\bf Z}$ 后，$e\beta$ 落在与权格 $\Lambda(H)$ 对偶的格 $\Lambda^*(H)$ 上；原文随即把这个对偶格解释为对偶群 $H^v$ 的权格。这一步依赖 $H$ 的表示与归一化约定（原文也明说 “factors of $e$ and $2\pi$ are convention dependent”），下面 $SO(3)$ 与 $SU(2)$ 的例子正是用 $T^3$ 的整/半整本征值来体现这种约定依赖。
\par
\par\medskip\noindent\hangindent=3.2em\hangafter=1 \textbf{[38]}~\textbf{1.9}~原文笔误，见原文 PDF 第 19 页。本页有三处拼写错误：“weekly coupled”应为“weakly coupled”，“monoples”应为“monopoles”，“superysymmetry”应为“supersymmetry”；另有“in the in the $H^v$-theory”重复了“in the”（在原文 PDF 第 19 页下半页）。物理上注意：Montonen–Olive 猜想交换的是微扰谱中的基本粒子与孤子，即 $H$ 理论与 $H^v$ 理论互换，耦合常数 $q/\hbar\leftrightarrow g/\hbar$；在自对偶的 $H=U(1)=H^v$ 情形，这一互换表现为 $e\leftrightarrow g=4\pi/e$，即强耦合与弱耦合互换。
\par
\par\medskip\noindent\hangindent=3.2em\hangafter=1 \textbf{[39]}~\textbf{1.10}~原文笔误，见原文 PDF 第 20 页：“half-odd inegral values”中 “inegral” 应为 “integral”；译文按正确含义处理。
\par
\par\medskip\noindent\hangindent=3.2em\hangafter=1 \textbf{[40]}~\textbf{1.10}~易错预警，见原文 PDF 第 20–21 页。核对 $\theta$ 项的归一化：$F^a_{\mu\nu}\tilde F^{a\mu\nu}$ 是全导数，$\partial_\mu K^\mu=\frac12\epsilon^{\mu\nu\rho\sigma}F^a_{\mu\nu}F^a_{\rho\sigma}$，故式 (44) 中系数取 $\theta e^2/32\pi^2$ 才能保证 $\theta\to\theta+2\pi$ 是物理上等价的周期；若把 $\tilde F$ 的 $\frac12$ 漏掉，$\theta$ 的周期会差一个因子 2。
\par
\par\medskip\noindent\hangindent=3.2em\hangafter=1 \textbf{[41]}~\textbf{1.10}~原文笔误，见原文 PDF 第 21 页：“with eigenvalues $q$ and $g$, respectivly”中 “respectivly” 应为 “respectively”。另，这里的 $n$ 是 $N$ 的本征值（整数），与 §1.9 中刻画 $\pi_1$ 的整数 $N$ 无关。
\par
\par\medskip\noindent\hangindent=3.2em\hangafter=1 \textbf{[42]}~\textbf{1.10}~易错预警，见原文 PDF 第 21 页式 (45)。这里 $g=\frac{4\pi}{e}m$ 的 $m$ 与电荷公式中 $ne$ 的 $n$ 是彼此独立的整数。对 't Hooft–Polyakov 单极，$g=-4\pi/e$ 意味着 $m=-1$（原文只写了“$n=1$”），代回式 (45) 才得到 $q=e(1-\theta/2\pi)$；若误取 $m=1$，会得到 $q=e(1+\theta/2\pi)$。
\par
\par\medskip\noindent\hangindent=3.2em\hangafter=1 \textbf{[43]}~\textbf{1.10}~见原文 PDF 第 22 页式 (46)–(48)。$\tau=\theta/2\pi+4\pi i/e^2$ 的实部为 $\theta/2\pi$、虚部为 $4\pi/e^2$；式 (48) 的 BPS 界即“质量不小于其格点到原点距离”，这正是后文（§2 起）把 BPS 态、格点与单值变换联系起来的出发点，建议牢记式 (46)–(48) 这三个式子。
\par
\par\medskip\noindent\hangindent=3.2em\hangafter=1 \textbf{[44]}~\textbf{2}~原文 PDF 第 22 页（印刷页 21）。本节从头建立 $N=1$ 超对称的旋量约定与超场形式，阅读时宜对照原文 PDF 第 22--23 页式 (49) 前后的约定表；文中所引 [22,23,24] 即 \cite{WB,West,Sohnius}（Wess--Bagger、West、Sohnius 三本/篇标准文献）。
\par
\par\medskip\noindent\hangindent=3.2em\hangafter=1 \textbf{[45]}~\textbf{2.1}~原文 PDF 第 22 页（印刷页 21）首句 "We start by descrbing our conventions" 中 "descrbing" 为原文拼写错误，应为 "describing"。译文按正确词义译出，原文误拼在此注明，未改动原文。
\par
\par\medskip\noindent\hangindent=3.2em\hangafter=1 \textbf{[46]}~\textbf{2.1}~原文 PDF 第 22 页，式 (49) 之前。$M^T\sigma_2M=\sigma_2$ 是由 $\epsilon^{\alpha\beta}=M^\alpha_{~\gamma}\epsilon^{\gamma\delta}M^\delta_{~\beta}$ 与 $\sigma_2=i\epsilon$ 对照得到的矩阵形式；注意 $M_\alpha^{~\beta}$（行 $\alpha$、列 $\beta$）与转置 $(M^T)_\alpha^{~\beta}=M^\beta_{~\alpha}$ 的指标对应，以及 $(M^{-1})^\alpha_{~\beta}$ 的指标位置，稍有不慎就会把升降指标的变换写反。
\par
\par\medskip\noindent\hangindent=3.2em\hangafter=1 \textbf{[47]}~\textbf{2.1}~原文 PDF 第 23 页（印刷页 22），$\bar\sigma^\mu$ 的定义式。由 $\epsilon$ 张量两次升指标得 $(\bar\sigma^\mu)^{\dot\alpha\alpha}=-(\sigma^\mu)^{\alpha\dot\alpha}=(i\sigma_2)(\sigma^\mu)^T(i\sigma_2)^T$，数值上为 $(1,-\vec\sigma)$，即 $\bar\sigma^0=\sigma^0=1$、$\bar\sigma^i=-\sigma^i$。这一负号是 (+---) 号差下 $\sigma$ 与 $\bar\sigma$ 关系的关键，后文 $W^\alpha W_\alpha$ 中 $F\widetilde F$ 项的符号由此决定；切勿套用 $(-,+,+,+)$ 文献的约定。
\par
\par\medskip\noindent\hangindent=3.2em\hangafter=1 \textbf{[48]}~\textbf{2.1}~原文 PDF 第 23 页，$\sigma^{\mu\nu}$ 定义与旋量标量积。两处易错：(i) 因子 $1/4$（而非某些教材的 $1/2$）来自 $\sigma^{\mu\nu}=\frac14(\sigma^\mu\bar\sigma^\nu-\sigma^\nu\bar\sigma^\mu)$，且第一个式子的指标为 $(\sigma^{\mu\nu})_\alpha^{~\beta}$、第二个为 $(\bar\sigma^{\mu\nu})^{\dot\alpha}_{~\dot\beta}$，缩并时方向不同；(ii) $\psi\chi=\psi^\alpha\chi_\alpha=-\psi_\alpha\chi^\alpha=\chi^\alpha\psi_\alpha=\chi\psi$ 中两个负号分别来自 $\epsilon$ 升指标与 Grassmann 反交换，$({\psi\chi})^\dagger=\bar\chi\bar\psi$ 的次序也不可颠倒。后文所有 Yukawa 项与质量项的符号都依赖这组约定。
\par
\par\medskip\noindent\hangindent=3.2em\hangafter=1 \textbf{[49]}~\textbf{2.1}~原文 PDF 第 23 页，$\gamma_5$ 与螺旋度。在 (+---) 号差下对 Dirac 方程的解有 $J_3=-\frac12\gamma_5$，故 $\gamma_5=+1$ 对应负螺旋度（原文写作 $+1\Rightarrow -ve$ helicity）。这与 $(-,+,+,+)$ 文献中常见的"$\gamma_5=+1$ 即正螺旋度"相反，引用时务必核对号差约定。
\par
\par\medskip\noindent\hangindent=3.2em\hangafter=1 \textbf{[50]}~\textbf{2.2}~原文 PDF 第 23 页（印刷页 22），无质量表示的 $\{Q^I_\alpha,\bar Q_{\dot\alpha J}\}$ 矩阵。取 $P^\mu=M(1,0,0,1)$ 后 $\sigma^\mu P_\mu=\mathrm{diag}(P^0-P^3,\,P^0+P^3)=\mathrm{diag}(0,2M)$，乘以式 (50) 的因子 2 即得右下角的 $4M$。又：由 $\{Q_1,\bar Q_{\dot1}\}=0$ 推出 $Q_1|phys>=0$ 的依据是幺正性，$0=\langle phys|\{Q_1,\bar Q_{\dot1}\}|phys\rangle=\|Q_1|phys\rangle\|^2+\|\bar Q_{\dot1}|phys\rangle\|^2$；注意超荷反对易子为 $+2\sigma^\mu P_\mu$（正号），而 §2.4 超协变导数为 $\{D_\alpha,\bar D_{\dot\alpha}\}=-2i\sigma^\mu\partial_\mu$（负号），两者恰好相反，是本章最常见的符号陷阱。
\par
\par\medskip\noindent\hangindent=3.2em\hangafter=1 \textbf{[51]}~\textbf{2.2}~原文 PDF 第 23 页。"$a^I$ is a rising operator and $(a^I)^\dagger$ is a lowering operator for the helicity" 一句容易误解：按角动量，$(a^I)^\dagger$ 使 $J_3$ 增加 $1/2$；但原文以"$m$ 个 $(a^I)^\dagger$ 给出螺旋度 $\lambda-m/2$"定义螺旋度，故对螺旋度而言 $(a^I)^\dagger$ 是下降算符。判断升降时请以 $J_3$ 的矩阵元为准，不要只看名称。
\par
\par\medskip\noindent\hangindent=3.2em\hangafter=1 \textbf{[52]}~\textbf{2.2}~原文 PDF 第 24 页（印刷页 23）。"we have to suplement them with their CPT conjugates" 与 "which can be assiged to them" 中 "suplement"、"assiged" 均为原文拼写错误，应作 "supplement"、"assigned"。又：表中 $N=2,\ \lambda=1/2$ 一行把 $|1/2>,2|0>,|-1/2>$ 与其 CPT 共轭 $|-1/2>,2|0>,|1/2>$ 并列，故 $|0>$、$|1/2>$ 各出现两次，这不是重复排版，而是把表示扩充为 CPT 自共轭的结果。
\par
\par\medskip\noindent\hangindent=3.2em\hangafter=1 \textbf{[53]}~\textbf{2.2}~原文 PDF 第 25 页（印刷页 24）。三处提醒：(i) 最大自旋为 $N/2$ 而非 $N$，关键是 $(a^I_1)^\dagger(a^I_2)^\dagger$ 在自旋指标上反对称、只能给出标量（原文用 $\frac12\epsilon^{\alpha\beta}(a^I_\alpha)^\dagger(a^I_\beta)^\dagger$ 说明）；(ii) $USp(2N)$ 为紧致辛群，有质量多重态的简并度按其不可约表示分类，原文明确略去不证；(iii) "the dimension of the the representation" 中衍出一个 "the"，系原文笔误。
\par
\par\medskip\noindent\hangindent=3.2em\hangafter=1 \textbf{[54]}~\textbf{2.3}~原文 PDF 第 26 页（印刷页 25），式 (51) 前后。两处易错与一处笔误：(i) 中心荷只出现在 $\{Q,Q\}$ 与 $\{\bar Q,\bar Q\}$ 中，系数为 $2\sqrt2\,\epsilon_{\alpha\beta}Z^{IJ}$（不是 $2\epsilon Z$），$\{Q,\bar Q\}$ 中没有中心荷；(ii) 斜对角化后 $Z=\epsilon\otimes D$，指标 $(a,m)$ 的拆分意味着不同 $m$ 的子代数相互对易，"只需考虑 $N=2$"的根据在此——每个 $m$ 给出独立的 $N=2$ 代数，各自给出 $M\geq\sqrt2|Z_m|$，故合起来要求 $M\geq\sqrt2\max_m|Z_m|$；(iii) 原文 "Haag, Lapuszanski and Sohnius" 中 "Lapuszanski" 应作 "Łopuszański"（即 Haag--Łopuszański--Sohnius 定理，文献 \cite{HLS}）。
\par
\par\medskip\noindent\hangindent=3.2em\hangafter=1 \textbf{[55]}~\textbf{2.3}~原文 PDF 第 26 页（印刷页 25），式 (53) 及其前一句。两点：(i) 原文写 "Then, the algebra (51) reduces to"，但按上下文此处应指斜对角化后得到的 $N=2$ 代数式 (52)（式 (51) 是带 $I,J$ 指标的完整代数），疑为原文式号笔误，译文照原文译出，请教授核对；(ii) 组合 $a_\alpha=\frac12\{Q^1_\alpha+\epsilon_{\alpha\beta}(Q^2_\beta)^\dagger\}$ 中的 $(Q^2_\beta)^\dagger$ 是 $\bar Q_{\dot\beta}^{\,2}$，其指标结构（带点指标 $\dot\beta$ 与 $a_\alpha$ 的 $\alpha$ 通过 $\epsilon$ 相配）是得到 $\delta_{\alpha\beta}(M\pm\sqrt2Z)$ 的关键，符号写反则长短多重态的判据会颠倒。
\par
\par\medskip\noindent\hangindent=3.2em\hangafter=1 \textbf{[56]}~\textbf{2.3}~原文 PDF 第 27 页（印刷页 26），式 (54)。这里 $a=|\langle\phi\rangle|$ 是 §1 中希格斯真空的模（原文 PDF 第 21 页明确取 $|\phi|=a$），$e$ 为基本电荷，$n,m$ 分别为电荷与磁荷的整数量子数，$\tau=\theta/2\pi+4\pi i/e^2$（原文 PDF 第 22 页式 (47)）；式 (54) 把中心荷归一化为 $ae(n+m\tau)$，故 $M\geq\sqrt2|Z|$ 与 Bogomol'nyi 界 (35) 同义。"量子修正不应产生额外自由度"一句是原文的启发式论证（短多重态的自由度数由代数决定，不能连续改变），并非严格证明。
\par
\par\medskip\noindent\hangindent=3.2em\hangafter=1 \textbf{[57]}~\textbf{2.4}~原文 PDF 第 27 页（印刷页 26），超空间约定。注意 $\theta\theta\equiv\theta^\alpha\theta_\alpha=-2\theta^1\theta^2$ 与 $\bar\theta\bar\theta\equiv\bar\theta_{\dot\alpha}\bar\theta^{\dot\alpha}=+2\bar\theta_{\dot1}\bar\theta_{\dot2}$ 的符号相反：这源于 $\theta$ 与 $\bar\theta$ 指标升降所用的同一 $\epsilon$ 张量的缩并方式；换用其他教材的 $\epsilon$ 约定时，凡 $\theta\theta$ 型的项都会整体变号。第三式 $\theta\sigma^\mu\bar\theta\,\theta\sigma^\nu\bar\theta=\frac12\theta\theta\bar\theta\bar\theta\,\eta^{\mu\nu}$ 是 Fierz 重排的基本恒等式，其系数 $1/2$ 与 $\eta^{\mu\nu}$ 的位置都不可改动。
\par
\par\medskip\noindent\hangindent=3.2em\hangafter=1 \textbf{[58]}~\textbf{2.4}~原文 PDF 第 28 页（印刷页 27），式 (57) 有一处原文公式笔误：$D_\alpha=\partial/\partial\theta^\alpha+i\sigma^\mu_{\alpha\dot\alpha}\partial_\mu$ 中漏写了 $\bar\theta^{\dot\alpha}$ 因子，正确形式应为 $D_\alpha=\partial/\partial\theta^\alpha+i\sigma^\mu_{\alpha\dot\alpha}\bar\theta^{\dot\alpha}\partial_\mu$；否则 $\{D_\alpha,\bar D_{\dot\alpha}\}$ 只能得到 $-i\sigma^\mu_{\alpha\dot\alpha}\partial_\mu$ 项（只有正确形式才给出原文所写的 $-2i\sigma^\mu\partial_\mu$）。正文公式按原文照录，正确形式在此注明。另需注意 $\{Q,\bar Q\}=+2i\sigma^\mu\partial_\mu$ 与 $\{D,\bar D\}=-2i\sigma^\mu\partial_\mu$ 差一负号，且 $Q_\alpha$ 与 $\bar Q_{\dot\alpha}$ 的一阶项符号相反（分别对应式 (55) 中 $x^\mu$ 位移的两项），这三处合起来是本章最易出错的地方。
\par
\par\medskip\noindent\hangindent=3.2em\hangafter=1 \textbf{[59]}~\textbf{2.4}~原文 PDF 第 28 页，式 (58) 的超场展开。系数 $1$（$f$）、$1$（$\theta\phi$）、$1$（$\theta^2m$）、$1/2$ 量的 $\theta^2\bar\theta^2d$ 都是相对于 $\theta\theta=\theta^\alpha\theta_\alpha$ 的约定；注意展开中 $\theta^2$ 项、$\bar\theta^2$ 项彼此独立，$\theta\sigma^\mu\bar\theta v_\mu$ 是唯一同时含 $\theta$ 与 $\bar\theta$ 的双线性项。$d(x)$ 的变分是全导数这一事实，是"超对称拉氏量 = 超场最高分量"这一构造原则的根据。
\par
\par\medskip\noindent\hangindent=3.2em\hangafter=1 \textbf{[60]}~\textbf{2.4}~原文 PDF 第 28 页（印刷页 27），式 (59) 及其下方。原文 "Here, $A$ and $\psi$ are the fermionic and scalar components respectively" 中 "fermionic" 与 "scalar" 的词序写反了：由式 (59) 可见 $A$ 是标量分量、$\psi$ 是费米分量。译文按正确物理含义译出，原文错误在此注明。另外，式 (59) 中 $A,\psi,F$ 都依赖 $y=x+i\theta\sigma\bar\theta$ 而非 $x$，这正是手征超场的本质（$\bar D\Phi=0$ 的解全部集中在 $(y,\theta)$ 上）；式 (62)--(63) 换回 $(x,\theta,\bar\theta)$ 后出现的 $\Box A$、$\partial_\mu\psi$ 等项正是 $y$ 与 $x$ 之差的高阶效应，作分量计算时不可遗漏。
\par
\par\medskip\noindent\hangindent=3.2em\hangafter=1 \textbf{[61]}~\textbf{2.4}~原文 PDF 第 29 页（印刷页 28），矢量超场与 Wess--Zumino 规范。四点是核对重点：(i) $V$ 的展开系数 $i/2$（$\theta^2$ 项）、$-i/2$（$\bar\theta^2$ 项）、负号（$\theta\sigma^\mu\bar\theta A_\mu$ 项）由 $V=V^\dagger$ 的实性条件固定，改动任一系数都会破坏实性；(ii) 规范变换为 $V\to V+\Lambda+\Lambda^\dagger$ 且 $\Lambda$ 必须是手征超场，这是保持 $V$ 为实超场且 $W_\alpha$ 手征的前提；(iii) 只有在这一规范下才有 $V^3=0$ 与 $V^2=\frac12A_\mu A^\mu\theta^2\bar\theta^2$，超场多项式才能截断；(iv) 规范固定条件 $C=M=N=\chi=0$ 与超对称不相容（WZ 规范破坏超对称但不破坏规范对称性），故按此规范作超对称变换时必须配合补偿变换；原文明确说明这一点，切勿把 WZ 规范下的分量变换直接当作超对称变换。
\par
\par\medskip\noindent\hangindent=3.2em\hangafter=1 \textbf{[62]}~\textbf{2.4}~原文 PDF 第 29--30 页（印刷页 28--29），场强的超场定义。注意阿贝尔情形 $W_\alpha=-\frac14\bar D^2D_\alpha V$（因子 $1/4$），非阿贝尔情形 $W_\alpha=\frac18\bar D^2 e^{2V}D_\alpha e^{-2V}$（因子 $1/8$）：两者归一化不同，比较系数时务必分清。非阿贝尔规范变换为 $e^{-2V}\to e^{-i\Lambda^\dagger}e^{-2V}e^{i\Lambda}$（左乘 $\Lambda^\dagger$、右乘 $\Lambda$），与常见的 $e^{2V}$ 写法的符号约定相反；式 (65) 中 $f^{abc}A^b_\mu A^c_\nu$ 与 $D_\mu\bar\lambda^a=(\partial_\mu+f^{abc}A^b_\mu)\bar\lambda^c$ 的符号由 $T^a$ 的伴随表示约定固定，换用 $T^a_{bc}=+if^{abc}$ 的教材时整体反号，需逐项核对。
\par
\par\medskip\noindent\hangindent=3.2em\hangafter=1 \textbf{[63]}~\textbf{2.5}~原文 PDF 第 30 页（印刷页 29），$\Phi_i^\dagger\Phi_j$ 的最高分量与自由拉格朗日量。核对要点：(i) $-\frac14A_i^\dagger\Box A_j-\frac14\Box A_i^\dagger A_j$ 两项经分部积分合成 $+\frac12\partial_\mu A_i^\dagger\partial^\mu A_j$（实标量动能正确）；(ii) 费米项 $-\frac i2\psi_j\sigma^\mu\partial_\mu\bar\psi_i+\frac i2\partial_\mu\psi_j\sigma^\mu\bar\psi_i$ 分部积分后化为 $-i\bar\psi_i\bar\sigma^\mu\partial_\mu\psi_i$，相对符号极易写反，可用厄米性校验；(iii) $F_i^\dagger F_j$ 为辅助场项，原文 "summing over $i=j$" 系笔误，意思应为"取 $i=j$ 并对 $i$ 求和"。又：原文此处把 $g^{ij}$ 记作场空间度规（通常教材写作 $g_{i\bar\jmath}$），本译本按原记号照录。
\par
\par\medskip\noindent\hangindent=3.2em\hangafter=1 \textbf{[64]}~\textbf{2.5}~原文 PDF 第 31 页（印刷页 30），$R$ 对称性。注意 $R$ 作用为 $\Phi'(x,\theta)=e^{2in\alpha}\Phi(x,e^{-i\alpha}\theta)$，故等价地 $\theta\to e^{+i\alpha}\theta$、$d^2\theta\to e^{-2i\alpha}d^2\theta$；于是 $\int d^2\theta\,{\cal W}$ 不变要求 ${\cal W}$ 中每一项的 $R$ 特征标之和为 $1$，而 $K$ 须为 $R$ 中性。原文此处 "Since $\theta\to e^{+i\alpha}\theta$, or $d^2\theta\to\ldots$, The R-character..." 中 "or" 应为 "i.e."、"The" 误作大写，均系原文笔误。术语上，本译本把 R-character 统一译作“$R$ 特征标”（即超场的 $R$ 荷）。$R$ 特征标的定义（$A\to e^{2in\alpha}A$，$\psi\to e^{2i(n-1/2)\alpha}\psi$，$F\to e^{2i(n-1)\alpha}F$）与 $d^2\theta$ 的权重 $e^{-2i\alpha}$ 必须配套，否则可重整性对 ${\cal W}$ 的约束会算错。
\par
\par\medskip\noindent\hangindent=3.2em\hangafter=1 \textbf{[65]}~\textbf{2.5}~原文 PDF 第 31 页（印刷页 30），$W^\alpha W_\alpha|_{\theta\theta}$ 与式 (66)--(67)。关键约定：$\widetilde F^{\mu\nu}\equiv\frac12\epsilon^{\mu\nu\rho\sigma}F_{\rho\sigma}$，故式中的 $\frac i4\epsilon^{\mu\nu\rho\sigma}F_{\mu\nu}F_{\rho\sigma}=\frac i2 F\widetilde F$；对偶项的归一化取决于此定义，若把 $\widetilde F$ 定义成无 $1/2$ 的形式，$F\widetilde F$ 的系数会差因子 2。$F\widetilde F$ 为全导数，在通常的超对称拉格朗日量中被丢弃，把它保留并配以 $\tau$ 的虚部（$\theta/2\pi$）即得式 (67)：其中 $-1/4g^2$ 来自 $\tau$ 实部 $4\pi/g^2$ 乘 $-\frac12$ 再乘 $1/8\pi$，$\theta/32\pi^2$ 来自 $\theta/2\pi$ 乘 $\frac12$ 再乘 $1/8\pi$，可用来校验系数。
\par
\par\medskip\noindent\hangindent=3.2em\hangafter=1 \textbf{[66]}~\textbf{2.5}~原文 PDF 第 32 页（印刷页 31），式 (68) 的排版笔误：第一项印作 $\int d^\theta$，缺少上标 $2$，应为 $\int d^2\theta\,W^\alpha W_\alpha$。正文公式按原文照录，正确形式在此注明；若照原文排版，$\int d^\theta$ 只是一个形式上缺了上标的写法，物理含义仍是 $\int d^2\theta$。
\par
\par\medskip\noindent\hangindent=3.2em\hangafter=1 \textbf{[67]}~\textbf{2.5}~原文 PDF 第 32 页（印刷页 31），式 (68) 下方。$N=1$ 超对称允许标量部分与杨--米尔斯部分取任意相对归一化（每项分别不变），原文把标量部分归一化为 $1$；下一节（原文 §2.6，见原文 PDF 第 32--33 页式 (69) 之后）将用 $\Phi\to\Phi/g$ 重新标度，使两个费米子的动能项归一化相同，这正是 $N=2$ 对 $N=1$ 拉格朗日量施加的第一个约束。这是理解 $N=2$ 构造的关键一步。
\par
\par\medskip\noindent\hangindent=3.2em\hangafter=1 \textbf{[68]}~\textbf{2.5}~原文 PDF 第 32 页（印刷页 31），式 (70) 的符号需核对。由式 (69) 的辅助场项 $+\frac{1}{2g^2}D^aD^a-D^aA^\dagger T^aA$ 消去 $D^a$（得 $D^a=g^2A^\dagger T^aA$）后，得到的应是 $+\frac12g^2(A^\dagger T^aA)^2$：原文式 (70) 第二项写作负号，会使标量势沿 $D$ 方向无下界，且与下一节式 (71) 的正定结果 $-\frac{1}{2g^2}{\rm Tr}([A^\dagger,A]^2)\geq0$ 不一致，疑为原文笔误。译文照录原文的负号，请教授核对原文 PDF 第 32 页式 (70)。
\par
\par\medskip\noindent\hangindent=3.2em\hangafter=1 \textbf{[69]}~\textbf{2.6}~见原文 PDF 第 33 页式 (71)。这里的符号与因子极易出错。把辅助场 部分的拉格朗日量写全为 $\frac{1}{g^2}{\rm Tr}(\frac{1}{2}D^aD^a+D^a[A^\dagger,A]^a+F^{\dagger a}F^a)$， $D$ 的代数方程给出 $D^a=-[A^\dagger,A]^a$；代回后 $\frac{1}{2}D^2$ 项只剩 $-\frac{1}{2g^2}{\rm Tr}([A^\dagger,A]^2)$。由于 $[A^\dagger,A]$ 反厄米，这个量非负， 这正是标量势的正定性来源。注意：(i) 系数 $1/2$ 与 $1/g^2$ 同时出现，漏掉任何一个都会 使 (71) 与 (72) 中的 $-\frac{1}{2}[A^\dagger,A]^2$ 对不上；(ii) 原文把二次项写成 $\frac{1}{2}DD$（$D$ 的指标缩并省略），不要误读为两个不同的 $D$ 场。
\par
\par\medskip\noindent\hangindent=3.2em\hangafter=1 \textbf{[70]}~\textbf{2.6}~见原文 PDF 第 33 页式 (72)。(72) 的第一行是全纯规范耦合 $\tau=\theta/2\pi+4\pi i/g^2$ 的 $N=1$ 超空间写法；第二行才是分量形式，其 Yang--Mills 项系数为 $1/g^2$，标量动能项系数也是 $1/g^2$，两者的相对归一化正是上一段所说的 "两个费米子动能归一化相同" 所固定。与 §1 的 Yang--Mills--Higgs 拉格朗日量逐项对照时， 注意这里 $A$ 带有因子 $1/g$（见正文对 $\Phi\to\Phi/g$ 的标度），所以 Higgs 势写成 $-(1/2g^2){\rm Tr}([A^\dagger,A]^2)$ 而非 §1 中的 $\lambda(\phi^\dagger\phi-v^2)^2$ 形式。 另外 $\theta$ 项系数 $g^2\theta/32\pi^2$ 中的 $g^2$ 来自 $\tau$ 的虚部与 $F\widetilde F$ 的归一化，切勿与 $\theta/32\pi^2$ 混淆。
\par
\par\medskip\noindent\hangindent=3.2em\hangafter=1 \textbf{[71]}~\textbf{2.6}~见原文 PDF 第 34 页 $\Psi^{(3)}$ 的表达式。原文此处排版有笔误： $\exp$ 的方括号内 $2gV(\cdots$ 只开了一个圆括号而没有闭合，源码中 $\backslash$left[ $\cdots$ $\backslash$right] 也不配对；按上下文（以及 $\bar\theta\bar\theta$ 投影）应为 $\Phi^\dagger(\cdots)\exp[2gV(\cdots)]|_{\bar\theta\bar\theta}$。 本书译文照录原文的写法，未作修补。
\par
\par\medskip\noindent\hangindent=3.2em\hangafter=1 \textbf{[72]}~\textbf{2.6}~见原文 PDF 第 34 页式 (74) 及其后。两处需留意。 (i) 原文 (74) 第二行写的是 $(\Phi^\dagger e^{2gV})^a{\cal F}_a(\Phi)$，指数是 $+2gV$；而 (72) 中规范不变的凯勒项是 $\Phi^\dagger e^{-2V}\Phi$。若取 ${\cal F}=\frac12\tau\Psi^2$ 回到 (72)，按规范不变性这里通常应是 $\Phi^\dagger e^{-2gV}$ （$g$ 因 $\Phi\to\Phi/g$ 的标度而出现；同一理由下 (72) 的指数也应是 $-2gV$）。 也就是说 (72) 与 (74) 两处的符号及 $g$ 因子互不自洽，至少有一处是原文笔误。 译文照录原文，读者自行推导有效作用量时请以规范不变性为准核对文献 [27,28] （Grimm；Bilal）的约定。 (ii) 紧接着"读出凯勒势"的句子在原文中括号不配对：${\rm Im}(\Phi^{\dagger a}{\cal F}_a(\Phi)$ 少一个右括号，译文照录；另外度规 $g_{ab}={\rm Im}(\partial_a\partial_b{\cal F})$ 严格说是混合分量 $g_{a\bar b}$ 的写法（凯勒度规的非零分量是 $\partial_a\partial_{\bar b}K$）， 原文指标约定如此。
\par
\par\medskip\noindent\hangindent=3.2em\hangafter=1 \textbf{[73]}~\textbf{2.6}~见原文 PDF 第 34 页末段。这里的逻辑是：$N=2$ 超对称把低能有效作用量 完全装入一个全纯函数 ${\cal F}(\Psi)$（预势）之中，因此确定量子修正后的 ${\cal F}$ 就等于确定整个低能理论。注意 ${\cal F}$ 与 $N=1$ 的超势 ${\cal W}$ 是两个不同的 对象（前者是全纯规范动力学的一部分，后者给出 Yukawa 耦合与 F 项），全文中 "prepotential" 一律译作"预势"。可重整性只允许二次 ${\cal F}$，故低能有效理论 中 ${\cal F}$ 的非二次部分必定来自量子修正，确定它便是非微扰问题。
\par
\par\medskip\noindent\hangindent=3.2em\hangafter=1 \textbf{[74]}~\textbf{2.7}~见原文 PDF 第 34 页 §2.7 开头。这里的论证是 §2.3 关于含中心荷 代数表示（原文 PDF 第 26 页起）的结论的直接应用：无中心荷时，有质量的 $N=2$ 多重态必须是"长"多重态，其分量数比超多重态（短多重态）多；只有中心荷 取特定值（BPS 条件 $M=\sqrt{2}|Z|$）时，有质量表示才能保持"短"，才能装进 超多重态。因此裸质量 $m_i$ 必须同时作为中心荷出现——这正是 §2.9 要计算的内容。
\par
\par\medskip\noindent\hangindent=3.2em\hangafter=1 \textbf{[75]}~\textbf{2.7}~见原文 PDF 第 35 页式 (75)。这是本片最易写错的超势之一，请注意： (i) 超势为 $\sqrt{2}\widetilde Q_i\Phi Q_i+m_i\widetilde Q_iQ_i$，其中 $\Phi=\Phi^aT^a$ 取伴随表示，$Q_i$ 与 $\widetilde Q_i$ 分别取 $N_c$ 与 $\bar N_c$ 表示，指标缩并后 是规范单态。(ii) 系数 $\sqrt2$ 由 $N=2$ 超对称固定，不能与质量项的相对系数任意 改变；它与式 (78) 中物质超流里的 $i\sqrt2$ 因子同源。(iii) 当 $A$ 取真空期望值 $a$ 时，$\sqrt2\widetilde q\,\Phi q$ 与 $m\widetilde q q$ 一起给出 $N=2$ 不变 的质量项，其结构保证了 BPS 关系 $M=\sqrt2|Z|$（见 §2.9）。(iv) 原文把 $N_f$ 个 超多重态的质量都取为对角形式，未写味混合；一般情形可通过对 $Q_i$ 作 $U(N_f)$ 旋转 对角化。
\par
\par\medskip\noindent\hangindent=3.2em\hangafter=1 \textbf{[76]}~\textbf{2.7}~见原文 PDF 第 35 页式 (76)。两处易错：(i) $D^a$ 中 $q^\dagger_i\lambda^a q_i$ 与 $-\widetilde q_i\lambda^a\widetilde q_i^\dagger$ 两项符号相反，这是因为 $Q$ 与 $\widetilde Q$ 的 $U(1)$ 荷相反；若把第二项的符号写错，消去 $D$ 后的标量势就不再 半正定。(ii) 这里的 $\lambda^a$ 是基本表示的生成元矩阵（$N_c\times N_c$）， 不是伴随表示的结构常数 $f^{abc}$；$\sum_a$ 是对规范群所有生成元求和， 与式 (71) 中 ${\rm Tr}$ 的写法只是记号不同。(iii) 原文把 $D$ 项与超多重态质量项 一起极小化并引用文献 [29,30]；这一步的物理内容是寻找标量势的平坦方向， 其分析结果直接给出式 (77) 的矩阵条件。
\par
\par\medskip\noindent\hangindent=3.2em\hangafter=1 \textbf{[77]}~\textbf{2.7}~见原文 PDF 第 36 页式 (77) 之前。由 ${\rm Tr}[(qq^\dagger-\widetilde q^\dagger\widetilde q)\lambda^a]=0$ 推出 (77)，用的是 $\lambda^a$ 在不可约表示下的完备性（Schur 引理）：若矩阵 $M$ 与 所有生成元对易，则 $M\propto{\bf 1}$。这里 $M=qq^\dagger-\widetilde q^\dagger\widetilde q$ 实际上满足的是更强的条件 ${\rm Tr}(M\lambda^a)=0$（对所有 $a$），即 $M$ 正交于 所有无迹生成元，故 $M\propto{\bf 1}$。注意常数写成 $c^2{\bf 1}_{N_c}$（$c$ 可正可负， 写 $c^2$ 便于下面开方），这种"用平方记常数"的写法在原文中出现多次。
\par
\par\medskip\noindent\hangindent=3.2em\hangafter=1 \textbf{[78]}~\textbf{2.7}~见原文 PDF 第 36 页第 1 种情形。此段有两处原文笔误，译文照录： (i) "它有 $N-N_f$ 个零本征值"应为 $N_c-N_f$（同段其余各处均用 $N_c$）； (ii) "但 $N_c=N_f-1$ 时除外，那时它完全破缺"应为 $N_c=N_f+1$：因为破缺后的 未破缺群是 $SU(N_c-N_f)$，只有当 $N_c-N_f=1$ 时它是平凡群 $SU(1)$，规范对称性 才完全破缺。另外，"$2N_fN_c-N_f^2$ 个夸克超场变重"这个数目恰等于被破缺的规范 生成元数 $\dim SU(N_c)-\dim SU(N_c-N_f)$，即 Higgs 机制中每个被吃掉一个手征超场； 这与"剩下 $N_f^2$ 个无质量超场作为破缺整体对称性的 Goldstone 玻色子"是自洽的。
\par
\par\medskip\noindent\hangindent=3.2em\hangafter=1 \textbf{[79]}~\textbf{2.7}~见原文 PDF 第 36 页第 2 种情形。原文标题写作 "$N_f\leq N$"， 但按本节开头的分类 "$N_f<N_c$ 与 $N_f\geq N_c$"，此处应为 $N_f\geq N_c$； 同段末尾 "For $N_f>N$, there are surviving $R$-symmetries" 中的 $N$ 也应为 $N_c$（且按字面应是 $N_f\geq N_c$ 的情形）。两处均疑为原文把 $N_c$ 简写成 $N$ 造成的笔误，译文照录原文。$\tilde v_i^{(i)}=\sqrt{|v_i^{(i)}|^2-c^2}$ 直接由 式 (77) 的对角分量 $|v_i^{(i)}|^2-|\tilde v_i^{(i)}|^2=c^2$ 得到，注意这里取 的是 $|\cdot|^2$ 而非 $(\cdot)^2$。
\par
\par\medskip\noindent\hangindent=3.2em\hangafter=1 \textbf{[80]}~\textbf{2.8}~见原文 PDF 第 37 页 §2.8 开头。$M\ge\sqrt2|Z|$ 是 §2.3 讨论含中心荷 代数表示时得到的结论（原文 PDF 第 26--27 页，见分片 C）；本节要证明它与 §1 的 Bogomol'nyi 界是同一个界。注意 $\sqrt2$ 因子是 $N=2$ 代数归一化的产物，本节 末尾（原文 PDF 第 41 页）由 (84) 之前的 $\tau_{cl}$ 表达式得到的 $M\ge\sqrt2|Z|$ 必须与之一致，切勿中途丢掉或改成 $\sqrt2/2$。
\par
\par\medskip\noindent\hangindent=3.2em\hangafter=1 \textbf{[81]}~\textbf{2.8}~见原文 PDF 第 37--38 页。原文说 "a slightly lengthy computation yields" 便跳过了全部中间步骤，这是本节的第一个"含糊"处：实际做法是把 $\delta{\cal L}_D$ 中的各项分部积分，把不是全导数的部分用运动方程或 $\sigma$ 矩阵恒等式归并，最后得到四个全导数的组合。核对要点是式 (78) 中四个 系数依次为 $\sqrt2$、$i\sqrt2$、$\sqrt2$、$-i\sqrt2$，其中两个 $i$ 的符号由 $\bar\psi\bar\sigma^\mu\partial_\mu\psi$ 与 $\partial_\mu\bar\psi\bar\sigma^\mu\psi$ 的相对符号决定；写成超流时再把 $\epsilon$、$\bar\epsilon$ 吸收进去，所以 实际的 $S^\rho_\alpha$ 与这里差一个 $\epsilon$ 因子（见式 (79)）。
\par
\par\medskip\noindent\hangindent=3.2em\hangafter=1 \textbf{[82]}~\textbf{2.8}~见原文 PDF 第 38 页式 (78)。底本（作者 1996 年的 LaTeX 源码）把本式 写在 \textbackslash{}beq...\textbackslash{}eeq（即 equation 环境）中，却在式中带了 \textbackslash{}nonumber；当年的 LaTeX 2e 在 equation 环境中忽略 \textbackslash{}nonumber，故原文 PDF 该式仍有编号 (78)。在载入 amsmath 的 现代编译下，\textbackslash{}nonumber（=\textbackslash{}notag）会真正取消编号，从而使 (78) 及以后所有公式的编号 比原文少 1。为满足“公式编号与原文逐一对应”的要求，本译本删去了这一多余的 \textbackslash{}nonumber（公式的数学内容一字未动），该式仍编为 (78)。
\par
\par\medskip\noindent\hangindent=3.2em\hangafter=1 \textbf{[83]}~\textbf{2.8}~见原文 PDF 第 38 页式 (78)--(79)。这是本节最容易混淆的约定： (i) 原文把 $\epsilon,\bar\epsilon$ 吸收进 $S^\rho$，故 (79) 中 $S^\rho$ 是 "带参数的"超流；后面式 (80) 之前的 $S^{\rho}_{gauge}$、式 (82) 以及 §2.9 中的 $\epsilon^\alpha S^\mu_{(1)\alpha}$ 都沿用这一约定。计算超荷对易子时必须先 把参数提出，否则会多出 $1/\sqrt2$ 与 $i$ 因子。 (ii) (78) 中 $\bar{\cal W}$ 与 ${\cal W}$ 的导数项分别带 $\bar\psi$ 与 $\psi$， 其 $i\sqrt2$ 的符号相反，这是下面 (82) 中相应项的来源。
\par
\par\medskip\noindent\hangindent=3.2em\hangafter=1 \textbf{[84]}~\textbf{2.8}~见原文 PDF 第 38 页式 (80)--(81)。这两组变分是后面全部计算的输入， 极易抄错：(i) $\delta A^a_\mu$ 中 $\epsilon$ 与 $\bar\epsilon$ 的位置、以及 $-i$ 与 $+i$ 的分配；$\delta D^a$ 没有 $i$ 因子，且两项的 $\bar\sigma^\mu$ 在 不同的两侧。(ii) 式 (81) 第三行 $\delta\bar\psi$ 中 $\bar\epsilon$ 配 $F^\dagger$、 $\epsilon$ 配 $D_\mu A^\dagger$，与第二行恰好"交叉"，写成同侧就会破坏超对称。 (iii) $\delta F$ 最后一项 $-2iT^aA\bar\epsilon\bar\lambda^a$ 是为了抵消 (69) 中 $A^\dagger T^a\lambda^a\psi$ 项里 $\delta\psi$ 的一部分而必需的， 系数 2 与生成元的归一化有关；若漏掉此项，超流 (82) 将不守恒。
\par
\par\medskip\noindent\hangindent=3.2em\hangafter=1 \textbf{[85]}~\textbf{2.8}~见原文 PDF 第 39 页式 (82) 之前。原文那句 "the supercurrent is not just the sum of the above, plus the gauge-covariantized version of $S^\rho_{matter}$" 措辞含混（直译会读成 "不仅是上述之和再加协变化版本"），其本意是：超流并不等于"上面两项之和" 加上"物质超流的规范协变化"，还必须补上来自 $D^aA^\dagger T^aA$ 与 $iA^\dagger T^a\lambda^a\psi$ 的额外项，否则超流不守恒。译文按此理解处理， 读者对照 (82) 与 (78) 即可看出第三、四行就是协变化后的物质项，而第二行末的 $A^\dagger T^aA$ 项则是额外贡献。
\par
\par\medskip\noindent\hangindent=3.2em\hangafter=1 \textbf{[86]}~\textbf{2.8}~见原文 PDF 第 39 页。两个超流都整体乘了 $g^2$（左端写成 $g^2S$）， 因为下面用 $g^2S^\mu$ 与规范费米子的归一化配套；还原物理超流时别忘了除以 $g^2$。 第二个超流是 $SU(2)_R$ 旋转的结果：$\lambda_2\to-\lambda_1$ 使末两项整体变号， 同时变换参数换成 $\epsilon'$，这不是笔误。注意原文说这对应于 "$(\lambda,\psi)$ 二重态作 $SU(2)_R$ 变换"，其具体实现是 $(\lambda_1,\lambda_2)\to(\lambda_2,-\lambda_1)$，即 $SU(2)_R$ 的一个 $90^\circ$ 转动（$\epsilon_{12}$ 型），读者可自行验证它保持 $N=2$ 代数不变。
\par
\par\medskip\noindent\hangindent=3.2em\hangafter=1 \textbf{[87]}~\textbf{2.8}~见原文 PDF 第 39 页。这两个 $\sigma$ 矩阵恒等式的最后一项 $\pm i\epsilon^{abcd}$ 的符号完全取决于 $\epsilon^{0123}$ 的取值约定 （以及 $\eta={\rm diag}(1,-1,-1,-1)$ 的号差），原文在 §1.1 的约定中固定了这一点。 若采用相反的 $\epsilon^{0123}$ 约定，化简后会得到 $-iF+\widetilde F$ 之类的 不同组合，进而使中心荷的实部/虚部归属颠倒。这是本节最容易出现"整体符号错误" 的地方，请务必核对原始约定。
\par
\par\medskip\noindent\hangindent=3.2em\hangafter=1 \textbf{[88]}~\textbf{2.8}~见原文 PDF 第 39--40 页。这一步的指标位置极易写错： $\sigma_{\nu\alpha\dot\alpha}\bar\lambda^{\dot\alpha}$ 与 $(\sigma^\nu\bar\sigma^\mu\lambda)_\alpha$ 的旋量指标都作用在 $\alpha$ 上， 而 $\sigma^\mu_{\alpha\dot\alpha}\bar\lambda^{\dot\alpha}A^\dagger T^aA$ 一项来自 $\epsilon\sigma^\mu\bar\lambda A^\dagger T^aA$ 中去掉 $\epsilon$。 另外 $\bar\lambda^{\dot\alpha}=i(\sigma_y\lambda^\dagger)^{\dot\alpha}$ 中的 $i$ 与 $\sigma_y$ 是 $(\frac12,0)$ 与 $(0,\frac12)$ 表示互换的标准结果， 原文在括号里特别注明 $\bar\lambda_{\dot\beta}=\lambda^\dagger_\beta$ 供核对。
\par
\par\medskip\noindent\hangindent=3.2em\hangafter=1 \textbf{[89]}~\textbf{2.8}~见原文 PDF 第 40 页。这里原文把两处"显然"的步骤跳过了，实际用到： (i) 等时反对易关系 $\{\lambda^{\dagger a}_{1\gamma},\lambda^b_{1\lambda}\} \propto\delta^{ab}\delta_{\gamma\lambda}$（规范费米子按 $\lambda\to g\lambda$ 归一化后，这正是第一行系数 $1/g^4$ 变成第二行 $1/g^2$ 的原因）； (ii) $\lambda_1$ 与 $\lambda_2$ 同属一个 $N=2$ 多重态，彼此反对易，因此两组 费米子不能交叉缩并，只能各自成对出现，这正是第一行两项的来历； (iii) $\sigma^{j\,T}$ 的出现来自交换 $\alpha\leftrightarrow\beta$ 时转置 $\sigma^j$。最后的 $\sigma^i\sigma_y\sigma^{j\,T}$ 只有其反对称部分存活。
\par
\par\medskip\noindent\hangindent=3.2em\hangafter=1 \textbf{[90]}~\textbf{2.8}~见原文 PDF 第 40--41 页。两点提示：(i) 由 $-2(\sigma_y)_{\alpha\beta}\delta^{ij}$ 写成 $2i\epsilon_{\alpha\beta}\delta^{ij}$ 用了 $(\sigma_y)_{\alpha\beta}=-i\epsilon_{\alpha\beta}$ 的约定（即 $\epsilon_{12}=+1$、$\sigma_y=\begin{pmatrix}0&-i\\ i&0\end{pmatrix}$）； 若把 $\epsilon_{\alpha\beta}$ 的符号约定弄反，中心荷的实部与虚部就会对调。 (ii) 从 $F^{a0i}D_iA^{\dagger a}$ 到全导数 $\partial_i(F^{a0i}A^{\dagger a})$ 需要：对磁部分用 Bianchi 恒等式 $D_i\widetilde F^{a0i}=0$，对电部分在源外区域 用运动方程 $D_iF^{a0i}=0$。原文只说 "one can easily show"，这一步实质上是把 $\partial_i(F A)$ 展开后丢掉 $D_iF$ 项；它正是 §1 中 Bogomol'nyi 表面项论证 的超对称版本。
\par
\par\medskip\noindent\hangindent=3.2em\hangafter=1 \textbf{[91]}~\textbf{2.8}~见原文 PDF 第 41 页电荷与磁荷的表达式。这是全节最易出归一化错误 的地方：(i) $Q_{ele}$ 与 $Q_{mag}$ 的分母都是 $ag$，但右端分别是 $gn_e$ 与 $\frac{4\pi}{g}n_m$；磁荷的 $4\pi/g$ 来自 't Hooft 张量 $\widetilde F$ 的归一化 （见 §1.4 关于电荷量子化的讨论，原文 PDF 第 7 页起）。(ii) $a$ 是 Higgs 真空中 $A$ 的取值，其归一化取决于生成元的选取（原文 $T^a$ 取基本表示， ${\rm Tr}(T^aT^b)=\frac12\delta^{ab}$）；若换成 $A=\frac a2\sigma_3$ 的写法， $Q_{ele}$ 的系数就会差因子 2。(iii) 两个公式都隐含了 $a$ 与 $g$ 的乘积 $ag$，这正是狄拉克量子化条件 $eg=2\pi$ 型结构的体现。
\par
\par\medskip\noindent\hangindent=3.2em\hangafter=1 \textbf{[92]}~\textbf{2.8}~见原文 PDF 第 41 页。此处 $Z=-ia(n_e+(4\pi i/g^2)n_m)$ 与下面 含 $\theta$ 的 $Z=a(n_e+\tau_{cl}n_m)$ 相差一个整体因子 $-i$（原文自己也说 "ignoring an overall factor of $-i$"）。两种写法不可混用：$|Z|$ 不受整体相位 影响，但 BPS 等号条件 $M=\sqrt2|Z|$ 中 $Z$ 的相位与超对称投影（哪一个超荷 被保留）有关，写有效作用量时一般用后面 $Z=an_e+a_Dn_m$ 的全纯形式。 另外此处的比较对象是式 (52)（原文 PDF 第 22 页附近，分片 C 的范围内）， 原文说 "comparing with (52)" 即指 §2.2 中中心扩张的代数形式。
\par
\par\medskip\noindent\hangindent=3.2em\hangafter=1 \textbf{[93]}~\textbf{2.8}~见原文 PDF 第 41 页。(i) 原文说 "the Witten effect described in section (1.10)"，但 §1 只有九个小节，"$\theta$ 参数与磁单极电荷" 是 §1.9 （原文 PDF 第 20 页）；此处 "(1.10)" 疑为原文编号笔误。 (ii) 复耦合 $\tau_{cl}=\theta/2\pi+4\pi i/g^2$ 的下标 $cl$ 表示这是经典值 （未含量子修正）；它与 (72) 第一行的 $\tau$ 一致，区别只在本书后文会用 量子修正后的 $\tau$。 (iii) $Q_{ele}$ 的移动量 $(\theta g^2/8\pi^2)Q_{mag}$ 与 $Q_{mag}=4\pi n_m/g$ 合起来给出 $Q_{ele}=g(n_e+\frac{\theta}{2\pi}n_m)$，这正是 $\tau_{cl}$ 的 实部来源；核对时注意 $\theta$ 的周期 $2\pi$。
\par
\par\medskip\noindent\hangindent=3.2em\hangafter=1 \textbf{[94]}~\textbf{2.8}~见原文 PDF 第 41 页末段。(i) 原文用对偶论证而不是直接计算来"导出" $Z=an_e+a_Dn_m$，这属于典型的"含糊点拨"处：严格证明需要 §3 的 $SL(2,Z)$ 变换性质（原文 PDF 第 43 页起，分片 E），这里只需记住 $(a,a_D)$ 在 $SL(2,Z)$ 下按二重态变换、且 $a_D=\partial{\cal F}/\partial a$。 (ii) 这里的 $a$ 是标量场 $A$ 的真空期望值（在 $SU(2)$ 中 $a$ 是 Cartan 子代数分量的值），$a_D$ 是"对偶"标量；$Z$ 用 $a$ 与 $a_D$ 表达是 Seiberg--Witten 理论的核心公式之一。 (iii) 注意 $\tau$ 与预势的关系：$\tau=\partial^2{\cal F}/\partial a^2$， 因此 $a_D=\partial{\cal F}/\partial a$ 与 $\tau$ 是同一预势的一阶、二阶导数， 后文 §3.6--3.7（原文 PDF 第 54、57 页）的单值变换与 $U(1)$ 的 $\beta$ 函数 都建立在这个关系上。
\par
\par\medskip\noindent\hangindent=3.2em\hangafter=1 \textbf{[95]}~\textbf{2.9}~见原文 PDF 第 42 页 §2.9 开头。原文写短多重态 "must satisfy the relation $M=\sqrt{2}Z$"，严格说应为 $M=\sqrt{2}|Z|$：$Z$ 含相位（§2.8 末尾 已指出其相位约定任意），短（BPS）多重态的质量只由 $|Z|$ 决定；原文此处的 $Z$ 被默认为正实数（在适当相位约定下）。译文照录原文。另注意此处逻辑： "$Z$ 不满足该关系" 是推出"中心荷还缺一块"的依据，而不是说前面的计算有错。
\par
\par\medskip\noindent\hangindent=3.2em\hangafter=1 \textbf{[96]}~\textbf{2.9}~见原文 PDF 第 42 页超势。请注意两项的相对系数：$\sqrt2\widetilde QQ\Phi$ 与 $m\widetilde QQ$ 必须同时出现且系数固定，否则 $N=2$ 会破缺为 $N=1$； 常见错误是只写质量项而忘记 $\sqrt2\widetilde Q\Phi Q$，或把 $\sqrt2$ 写成 2。 这与 §2.7 式 (75) 是同一个超势，只是这里写出了味指标求和并明确 $Q_i$ 在 $N$、$\widetilde Q_i$ 在 $\bar N$ 表示。
\par
\par\medskip\noindent\hangindent=3.2em\hangafter=1 \textbf{[97]}~\textbf{2.9}~见原文 PDF 第 42 页这两组式子。要点：(i) 替换规则 $q\to\widetilde q^\dagger$、$\widetilde q^\dagger\to-q$ 只作用于标量 （$q,\widetilde q^\dagger$ 构成 $SU(2)_R$ 二重态），费米子 $\psi_q,\psi_{\widetilde q}$ 是单态因而不变；这与 §2.8 中 $\lambda_1\to\lambda_2,\lambda_2\to-\lambda_1$ 是同一个 $SU(2)_R$ 旋转。 (ii) $m$ 项前的 $i\sqrt{2}$ 因子来自式 (78) 中 $i\sqrt{2}\,\partial\bar{\cal W}/\partial A^\dagger\,\bar\psi\bar\sigma^\rho\epsilon$ 一项；因为 $\partial\bar{\cal W}/\partial\bar q = m\widetilde q$ 型， 代入后即得这些项。(iii) 第二组式子中 $m$ 项的整体符号翻转，正是因为 $\widetilde q^\dagger\to-q$；若符号写错，反对易子中的附加项 $\sum_i m_i S_i$ 就会抵消而不是相加。
\par
\par\medskip\noindent\hangindent=3.2em\hangafter=1 \textbf{[98]}~\textbf{2.9}~见原文 PDF 第 42--43 页 $S_i$ 的定义。这是标准的 Noether 荷： 第一项 $D_0q^\dagger q+qD_0q^\dagger$ 是标量场的"动量"部分（对称化以保证 厄米性），中间两项是费米子双线性的自旋贡献，括号内最后的替换 $(q\to\widetilde q,\ \psi_q\to\psi_{\widetilde q})$ 表示对 $\widetilde q$ 取相反的 $U(1)$ 荷（即整项变号）。原文说 "$i$ is not summed over"， 即 $S_i$ 对应第 $i$ 个味，而非求和。注意 $\psi^\dagger_{qi}\psi_{qi}$ 与 $\psi_{qi}\psi^\dagger_{qi}$ 的次序在两项中相反，这是费米子算符次序所致， 保持原文即可。
\par
\par\medskip\noindent\hangindent=3.2em\hangafter=1 \textbf{[99]}~\textbf{2.9}~见原文 PDF 第 43 页式 (84)。这是本片的最后结果，也是 §2 的核心公式 之一。系数 $1/\sqrt2$ 的来源值得记住：对 BPS 态 $M=\sqrt2|Z|$，而超多重态的 质量项直接进入 $M$，故质量 $m_i$ 必须以 $m_i/\sqrt2$ 的形式出现在 $Z$ 中， 再由 $S_i=\pm1$（对 $U(1)$ 荷为 $\pm1$ 的态）给出 $M=m_i$。若把 $1/\sqrt2$ 写成 $\sqrt2$，BPS 谱就会整体差因子 2。另外式 (84) 是 §3、§4 中确定 Seiberg--Witten 曲线质量参数的基础：裸质量 $m_i$ 正是通过这一项与 $a,a_D$ 一起进入中心荷，从而决定 BPS 谱的谱线。
\par
\par\medskip\noindent\hangindent=3.2em\hangafter=1 \textbf{[100]}~\textbf{3}~上式中 $Z=a(n_e+\frac{i}{\alpha}n_m)$ 在 $n_e\leftrightarrow n_m$、$\alpha\leftrightarrow1/\alpha$、$a\leftrightarrow a/\alpha$ 下只差一个整体相位（$Z\to iZ$），故 $|Z|$ 不变；这正是 $g\leftrightarrow g_D=4\pi/g$ 的 Montonen--Olive 对偶的“弱形式”。注意这只是经典（树图）层面的观察：$N=2$ 量子理论中 $Z$ 的重整化通过 $a_D=\partial{\cal F}/\partial a$ 进入，精确公式是 (3.12)。
\par
\par\medskip\noindent\hangindent=3.2em\hangafter=1 \textbf{[101]}~\textbf{3}~$\tau=\theta/2\pi+i4\pi/g^2$ 的定义极易记错：虚部是 $4\pi/g^2$（不是 $1/g^2$，也不是 $g^2/4\pi$），因此 $\operatorname{Im}\tau=4\pi/g^2>0$ 恰对应弱耦合极限 $g^2\to0$，即 $\tau$ 位于上半平面（upper half plane）。与本节导语的 $\alpha=g^2/4\pi$ 对照有 $\tau=\theta/2\pi+i/\alpha$。还要区分：此处的 $\tau$ 是经典（树图）耦合；在量子模空间上 $\tau(a)=\partial^2{\cal F}/\partial a^2$（见式 (3.5)），两者仅在 $|a|\to\infty$ 的渐近自由区重合。
\par
\par\medskip\noindent\hangindent=3.2em\hangafter=1 \textbf{[102]}~\textbf{3.1}~原文 “containts” 系 “contains” 的拼写笔误（原文 PDF 第 43 页）。势能前的因子 $1/2g^2$ 与 $\operatorname{Tr}$ 的归一化约定绑定（见 §2 的约定）；若改用基本表示下的 $\operatorname{tr}$，前置因子会随之改变。$V\geq0$，且 $V=0\iff[\phi,\phi^\dagger]=0$，这正是下文定义希格斯真空的依据。
\par
\par\medskip\noindent\hangindent=3.2em\hangafter=1 \textbf{[103]}~\textbf{3.1}~这里的逻辑链值得补全：$G/H$ 中保持嘉当子代数不变的元素构成 Weyl 群，它通过共轭作用把 $\phi$ 在嘉当内的坐标作反射；物理真空由 Weyl 轨道分类，故模空间的坐标函数必须是 Weyl 不变量。Chevalley 定理进一步保证：Weyl 不变多项式环是自由多项式环，恰有 $r$ 个生成元，这与“模空间维数 = 秩 $r$”一致。原文只说“Weyl 反射”而未点明这层不变量理论，容易让读者误以为 $G/H$ 全是规范冗余。
\par
\par\medskip\noindent\hangindent=3.2em\hangafter=1 \textbf{[104]}~\textbf{3.1}~“容易验证”一步的细节（原文 PDF 第 44 页）：$\phi={1\over2}a\sigma_3$ 时 $\phi^2={1\over4}a^2\mathbb{1}$，在基本表示中 $\operatorname{Tr}\mathbb{1}=2$，故 ${\rm Tr}(\phi^2)={1\over2}a^2$；Weyl 反射 $a\to-a$ 使之不变，而秩为 1 说明这是唯一的独立不变量。
\par
\par\medskip\noindent\hangindent=3.2em\hangafter=1 \textbf{[105]}~\textbf{3.1}~$({\rm Tr}\phi)^2=0$、$({\rm Tr}\phi)^3=0$ 是缩写，应理解为把 ${\rm Tr}\phi=\sum_i a_i$ 代入后恒为零所得的牛顿恒等式：由 $(\sum_i a_i)^2=0$ 得 $\sum_{i<j}a_ia_j=-{1\over2}\sum_i a_i^2$，于是 $u={1\over2}\sum_i a_i^2=-\sum_{i<j}a_ia_j$；由 $(\sum_i a_i)^3=0$ 并利用三变量对称性可得 $a_1a_2a_3={1\over3}\sum_i a_i^3$。按字面读会把这两个恒等式误解为空话。
\par
\par\medskip\noindent\hangindent=3.2em\hangafter=1 \textbf{[106]}~\textbf{3.1}~原文此处有一个符号笔误（原文 PDF 第 44 页）：与上一页的一般展开 $\det(\lambda-\phi)=\lambda^N-\frac12{\rm Tr}(\phi^2)\lambda^{N-2}-\frac13{\rm Tr}(\phi^3)\lambda^{N-3}+\cdots$ 对照，$N=3$ 时常数项为 $-\frac13{\rm Tr}(\phi^3)=-(a_1a_2a_3)$，故正确的定义应是 $v=+\frac13{\rm Tr}(\phi^3)=a_1a_2a_3$；原文写成 $v=-\frac13{\rm Tr}(\phi^3)=a_1a_2a_3$，中间表达式的负号有误（右端 $a_1a_2a_3$ 才是 $v$ 的定义）。公式已按原文照录，此处如实指出。
\par
\par\medskip\noindent\hangindent=3.2em\hangafter=1 \textbf{[107]}~\textbf{3.1}~原文用 $<$ 与 $>$ 表示真空期望值（$u=<{\rm Tr}(\phi^2)>$），这是老式 LaTeX 写法，排版上与不等号同形；规范记法应为 $\langle{\rm Tr}(\phi^2)\rangle$。数学内容未改动，仅在此提示。
\par
\par\medskip\noindent\hangindent=3.2em\hangafter=1 \textbf{[108]}~\textbf{3.2}~这一步“容易翻译”的实质：超荷的微分表示中含 $\partial/\partial\theta$ 项，要使 $\{Q,\bar Q\}\sim\partial_\mu$ 与 $Q$ 作用于 $\theta$ 的形式不变，$\theta$ 必须与 $Q$ 反向变换；同时 $d^4\theta$ 测度在 $\theta\to e^{i\alpha}\theta$ 下不变，而 $d^2\theta$ 测度带荷 $-2$。这既解释了 (3.1) 中 $\Phi$ 前的因子 $e^{2i\alpha}$（$\Phi$ 的 $U(1)_{\cal R}$ 荷为 2），也解释了超势 ${\cal W}=\sqrt{2}\widetilde Q\Phi Q$ 的总荷必须为 2（见 (3.1) 之后关于 ${\cal W}$ 的讨论）。
\par
\par\medskip\noindent\hangindent=3.2em\hangafter=1 \textbf{[109]}~\textbf{3.2}~注意区分 $U(1)_J$ 与 $U(1)_{\cal R}$：在 (3.1) 中 $V$ 的变换完全相同，区别仅在于手征超场 $\Phi$ 是否带整体相位 $e^{2i\alpha}$。$U(1)_J\subset SU(2)_R$ 由 $\sigma_3$ 生成，是 $N=1$ 语言下仍显式的那个 $R$ 对称性；$U(1)_{\cal R}$ 则是 $N=2$ 的额外整体 $U(1)$。两者是独立参数，切勿把它们混成一个 $\alpha$（原文在 (3.2) 中同样用 $\alpha$ 记两个不同 $U(1)$ 的参数）。
\par
\par\medskip\noindent\hangindent=3.2em\hangafter=1 \textbf{[110]}~\textbf{3.2}~由 (3.1) 中 $\Phi\to e^{2i\alpha}\Phi$ 可知 ${\cal W}=\sqrt{2}\widetilde Q\Phi Q$ 的 $U(1)_{\cal R}$ 荷为 $\Phi$ 的荷 2；要使超势不变，$Q,\widetilde Q$ 必须中性——这正是 (3.2) 中它们在 $U(1)_{\cal R}$ 下只随超空间坐标平移而无整体相位的理由。
\par
\par\medskip\noindent\hangindent=3.2em\hangafter=1 \textbf{[111]}~\textbf{3.2}~$\psi_D$ 的 $U(1)_{\cal R}$ 变换 $e^{i\alpha\gamma_5}$ 是轴（手征）转动，会在路径积分测度中产生反常，这是 $U(1)_{\cal R}$ 破缺的根源；而 $U(1)_J$ 给出普通的矢量相位 $e^{i\alpha}$（$\lambda$ 与 $\psi$ 反向变换，使 $\psi_D=(\lambda,\bar\psi)$ 获得同一相位），因而是非反常的。两种 $U(1)$ 对手征反常的贡献截然不同，这是本节最易混淆之处。
\par
\par\medskip\noindent\hangindent=3.2em\hangafter=1 \textbf{[112]}~\textbf{3.2}~指标定理给出的零模数核对：对 $SU(N_c)$，表示 $R$ 中一个左手 Weyl 费米子在 $\nu$ 个瞬子背景中有 $2T(R)\nu$ 个零模；$T({\rm adj})=N_c$（故 $2N_c$ 个），$T({\rm fund})={1\over2}$（故 1 个），与原文一致。原文此处 “left moving”（左行）与 “left-handed”（左手）混用，实指同一件事（手征零模）；“representaion” 是 “representation” 的拼写笔误（原文 PDF 第 46 页）。用集体坐标语言说：瞬子振幅的费米子积分要恰好插入全部零模才非零，下式 (3.3) 就是能“吃掉”零模的最小关联函数，其 $U(1)_{\cal R}$ 荷决定破缺后的离散群。
\par
\par\medskip\noindent\hangindent=3.2em\hangafter=1 \textbf{[113]}~\textbf{3.2}~荷的核对（极易错）：由 (3.1) 下方的分量表，$\lambda,\psi$ 的 $U(1)_{\cal R}$ 荷各为 $+1$，$\psi_q,\widetilde\psi_q$ 各为 $-1$，故 (3.3) 的总荷为 $2N_c\cdot1+2N_c\cdot1-N_f\cdot1-N_f\cdot1=4N_c-2N_f$。令上式的相位为 1 得 $\alpha=2\pi n/(4N_c-2N_f)$。注意原文在离散情形把群元素写成 $e^{2\pi i\alpha}$、$\alpha=n/4N_c$，这与前文连续转动参数（荷 $q$ 的场取 $e^{iq\alpha}$）的归一化不同：离散情形下荷 $q$ 的场实际获得 $(e^{2\pi in/(4N_c)})^q$。下文说 $\phi^2$（荷 4）变换为 $e^{2\pi in/N_c}\phi^2$ 即按此约定，可自行核对。
\par
\par\medskip\noindent\hangindent=3.2em\hangafter=1 \textbf{[114]}~\textbf{3.2}~$n=2N_c$ 对应的群元素为 $e^{2\pi i\cdot2N_c/(4N_c)}=e^{i\pi}=-1$，作用在荷 1 的费米子上给 $-1$，恰与 $SU(2)_R$ 中心 $\lambda,\psi\to e^{i\pi}(\lambda,\psi)$ 相同，故必须取商 ${\bf Z}_2$。这是下文 $SU(2)$ 最终群 $(SU(2)_R\times{\bf Z}_4)/{\bf Z}_2$ 的来源。
\par
\par\medskip\noindent\hangindent=3.2em\hangafter=1 \textbf{[115]}~\textbf{3.2}~“$\phi^3$ 的不变性挑出 $n=2N_c,4N_c$”是压缩表述（原文 PDF 第 47 页），容易误读：单看 $\phi^3$ 的不变性（相位 $e^{3\pi i n/N_c}=1$）只给出 $n$ 为偶数，即 6 个元素（${\bf Z}_6$，与下一句“破缺掉的 ${\bf Z}_6$”呼应）；必须与 $\phi^2$ 的条件 $n\in\{N_c,2N_c,3N_c,4N_c\}=\{3,6,9,12\}$ 取交集，才剩下 $n=6,12$，即 ${\bf Z}_2$。
\par
\par\medskip\noindent\hangindent=3.2em\hangafter=1 \textbf{[116]}~\textbf{3.2}~对 $SU(4)$ 及以上“没有任何子群存活”一句若按字面理解会与计算矛盾（原文 PDF 第 47 页）：元素 $n=2N_c$ 对所有 Weyl 不变量 ${\rm Tr}\phi^k$（$U(1)_{\cal R}$ 荷 $q=2k$，均为偶数）给出的相位为 $e^{2\pi i q\cdot2N_c/(4N_c)}=e^{i\pi q}=1$，即它保持整个真空模空间不变。该元素恰好是 $SU(2)_R$ 的中心（见上一段），在 $(SU(2)_R\times{\bf Z}_{4N_c})/{\bf Z}_2$ 的商下与群中心重合，因而不给出新的对称性。原文的意思应理解为“商掉中心后不再有新的离散 $R$ 对称性”；这一点原文没有点明，特此提醒。
\par
\par\medskip\noindent\hangindent=3.2em\hangafter=1 \textbf{[117]}~\textbf{3.3}~原文 “non of the massive states” 应为 “none of the massive states”；“as we will in the remaining part of this section” 语句不完整（疑脱漏 “do/describe”）。两处均为原文笔误（原文 PDF 第 47 页），译文按语义补足，未增删物理内容。
\par
\par\medskip\noindent\hangindent=3.2em\hangafter=1 \textbf{[118]}~\textbf{3.3}~指标位置极易写反：$K={\rm Im}(\bar A_i\,\partial{\cal F}/\partial A_i)$，度规 $g_{i\bar j}={\rm Im}\,\partial^2{\cal F}/\partial a_i\partial a_j$；注意 $\bar A_i$ 是复共轭并带下指标。另注意归一化：由 (3.4) 的 $1/4\pi$，(3.5) 直接给出 $\tau=\partial^2{\cal F}/\partial a^2$（不再带 $1/4\pi$）；若漏掉 (3.4) 中的 $1/4\pi$，$\tau$ 会与 $\tau_{cl}=\theta/2\pi+i4\pi/g^2$ 相差 $4\pi$，后文所有对偶变换都会随之错位。
\par
\par\medskip\noindent\hangindent=3.2em\hangafter=1 \textbf{[119]}~\textbf{3.3}~“度规 = 广义规范耦合”是 $N=2$ 的核心结构：同一个全纯函数 ${\cal F}$ 同时给出标量场空间的凯勒度规与规范动力学因子（$\operatorname{Im}\tau\,F_{\mu\nu}F^{\mu\nu}$ 的系数）。由于 $N=2$ 的非重整化性质，这一等同关系不因量子修正而改变；这正是后文能用 ${\rm Im}\,\tau>0$（动能正定）约束 $\tau$、并把求度规归结为求 $a_D(u),a(u)$ 的原因。
\par
\par\medskip\noindent\hangindent=3.2em\hangafter=1 \textbf{[120]}~\textbf{3.3}~记号陷阱：此处的 $\alpha$ 是 $U(1)_{\cal R}$ 的转动参数；本节导语中的 $\alpha=g^2/4\pi$；(3.10) 前后 $SL(2,{\rm Z})$ 矩阵的四个元素也记作 $\alpha,\beta,\gamma,\delta$。同一个字母在本节中有三种完全不同的含义，阅读时必须按上下文区分，这是原文记号最易误导之处。
\par
\par\medskip\noindent\hangindent=3.2em\hangafter=1 \textbf{[121]}~\textbf{3.3}~反常系数的归一化易错：原文用 $(32\pi^2)^{-1}\int F\widetilde F\in{\bf Z}$（分母是 $32\pi^2$ 而非 $16\pi^2$），这与 $\theta$ 项的写法 $\theta\,(32\pi^2)^{-1}\int F\widetilde F$ 配套（可由 (3.4)--(3.5) 中 $\theta$ 项的系数验证）。核对离散子群时必须使用同一约定，并记住这里的 $\alpha$ 是 $U(1)_{\cal R}$ 的转角。
\par
\par\medskip\noindent\hangindent=3.2em\hangafter=1 \textbf{[122]}~\textbf{3.3}~$N_c\to N_c-N_f/2$ 的来源是零模计数：伴随表示（$\lambda,\psi$，荷各 $+1$）贡献 $2N_c+2N_c=4N_c$，每个基本表示味（$\psi_q,\widetilde\psi_q$，荷各 $-1$）贡献 2，故有效荷为 $4N_c-2N_f=4(N_c-N_f/2)$，离散群为 ${\bf Z}_{4(N_c-N_f/2)}={\bf Z}_{4N_c-2N_f}$。
\par
\par\medskip\noindent\hangindent=3.2em\hangafter=1 \textbf{[123]}~\textbf{3.3}~手征转动把 $\theta$ 移走的机制：$U(1)_{\cal R}$ 转动使 $\delta{\cal L}_{eff}=-\alpha N_c/(8\pi^2)F\widetilde F$，而 $\theta$ 项为 $\theta\,(32\pi^2)^{-1}F\widetilde F$，故等价于 $\theta\to\theta-4\pi N_c\alpha$（纯 Yang--Mills，$N_f=0$）；选取适当的 $\alpha$ 即可把 $\theta$ 归零。此后本节默认 $\theta=0$，从而 $\tau_{cl}=i4\pi/g^2$ 为纯虚数。
\par
\par\medskip\noindent\hangindent=3.2em\hangafter=1 \textbf{[124]}~\textbf{3.3}~“容易积分”的细节（原文 PDF 第 49 页）：由 $\partial^3{\cal F}/\partial A^3=(N_c/\pi)(i/A)$ 作两次不定积分得 $\partial^2{\cal F}/\partial A^2=(iN_c/\pi)\ln A+c_1$；$N_c=2$ 时与 $\operatorname{Im}\tau$ 的单圈形式 $\tau={\cal F}''=(i/\pi)(\ln(A^2/\Lambda^2)+3)$ 比较可见，$-\ln\Lambda^2$ 与常数 3 都属于积分常数（中含有限部分）。这些常数要与经典项 ${\cal F}=\frac12\tau_{cl}A^2$ 的归一化匹配才能唯一确定，原文省略了这一步；常数“3”也可直接对 $A^2\ln A^2$ 求二阶导（得 $2+1$）核对。
\par
\par\medskip\noindent\hangindent=3.2em\hangafter=1 \textbf{[125]}~\textbf{3.3}~“$N=2$ 的 $\beta$ 函数只有单圈”是 $N=2$ 非重整化定理的结论。系数可核对：在 $N=1$ 语言中单圈系数为 $3C_2(G)-T({\rm adj})$，对 $SU(N_c)$ 即 $3N_c-N_c=2N_c$，故 $\beta(g)=-g^3\cdot 2N_c/(16\pi^2)=-N_c g^3/8\pi^2$；$N_c=2$ 正给出原文的 $\beta(g)=-g^3/4\pi^2$。正因为如此，“反常严格单圈”与“预势在微扰论中精确到单圈”才能同时成立。对比 $N=1$：那里手征反常仍严格单圈，但 $\beta$ 函数会得到高阶修正（见 [36,37]），所以 $N=1$ 的预势没有这种单圈精确性。
\par
\par\medskip\noindent\hangindent=3.2em\hangafter=1 \textbf{[126]}~\textbf{3.3}~单圈 $\beta$ 函数积分给出 $8\pi^2/g^2(\mu)=4\ln(\mu/\Lambda)$：由 $dg/d\ln\mu=-g^3/4\pi^2$ 得 $d(8\pi^2/g^2)/d\ln\mu=4$，于是 $e^{-8\pi^2k/g^2}=(\Lambda/\mu)^{4k}$；取 $\mu\sim a$（相关能标）即得下式，常数因子吸收进 $\Lambda$ 的定义。另：全纯性论证（见 [38,39]）排除反瞬子修正——反瞬子贡献依赖 $\bar g$ 而不全纯。
\par
\par\medskip\noindent\hangindent=3.2em\hangafter=1 \textbf{[127]}~\textbf{3.3}~给 $\Lambda$ 赋 $U(1)_{\cal R}$ 荷 2 的技巧值得说明：$\Lambda$ 是动力学标度而非基本场，但在反常意义下可以赋荷，使得 $\ln(A^2/\Lambda^2)$ 成为中性（$A$ 的 $U(1)_{\cal R}$ 荷为 2，故 $A^2$ 与 $\Lambda^2$ 的荷都是 4）；于是单圈项整体带荷 4，与 $k$-瞬子项 $(\Lambda/A)^{4k}A^2$ 的荷一致，$U(1)_{\cal R}$ 便以齐次方式恢复。这是 Seiberg 的全纯性论证（见 [38,39]）的关键一步。
\par
\par\medskip\noindent\hangindent=3.2em\hangafter=1 \textbf{[128]}~\textbf{3.3}~“取 $A$ 为积分限”的细节：由 $\beta$ 函数积分 $1/g^2$ 得 $4\pi/g^2(a)=(2/\pi)\ln(a/\Lambda)+{\rm const}$，其虚部形式恰给出后文 $\operatorname{Im}\tau=\frac1\pi(\ln\frac{a^2}{\Lambda^2}+3)$ 的对数项；常数部分（含“3”）由与经典项 $\frac12\tau_{cl}A^2$ 的匹配确定。原文 “sompnent” 是 “component” 的拼写笔误（原文 PDF 第 50 页），“vev” 指 vacuum expectation value。
\par
\par\medskip\noindent\hangindent=3.2em\hangafter=1 \textbf{[129]}~\textbf{3.3}~“调和函数不可能有整体极小值”用的是调和函数的最小值原理（maximum principle）：非常数调和函数在开区域内部不能取极值。这里有两层意思：${\rm Im}\,\tau$ 必须处处为正（否则度规不定、动能项变负），但单圈表达式 ${\rm Im}\,\tau=\frac1\pi(\ln\frac{a^2}{\Lambda^2}+3)$ 在小 $|a|$ 处会变负，故微扰描述只在弱耦合区（大 $|a|$）有效；$\tau(a)$ 的多值性（$\ln a^2$ 的对数支）正是后文单值变换（monodromy）的几何根源（原文 PDF 第 50 页）。
\par
\par\medskip\noindent\hangindent=3.2em\hangafter=1 \textbf{[130]}~\textbf{3.4}~这里的 $8\pi$ 与 §1 的归一化约定严格绑定：把 $U(1)\subset SU(2)$ 归一化到“$SO(3)$ 场整数电荷”时，基本表示物质场的电荷为半整数，磁单极方程右端为 $8\pi\delta^{(3)}(x)$。若采用 $eg=2\pi$ 或 $eg=4\pi$ 等不同约定，$8\pi$、$1/8\pi$、$1/16\pi$ 等因子会整体改变；务必与 §1（1.2--1.4 小节）中电荷量子化与规范场归一化的约定对照核对（原文 PDF 第 51 页）。
\par
\par\medskip\noindent\hangindent=3.2em\hangafter=1 \textbf{[131]}~\textbf{3.4}~中间略去的代数可补出：利用 2-形式的恒等式 $\widetilde F_D F=F_D\widetilde F$、$\widetilde F_D\widetilde F=-F_D F$ 及 $\operatorname{Re}(iX)=-\operatorname{Im}X$，可把 $\frac1{16\pi}\operatorname{Re}\int(\widetilde F_D-iF_D)(F+i\widetilde F)$ 化为 $\frac1{8\pi}\int\widetilde F_D F$；若把 $X=F+i\widetilde F$ 当作整体（注意 $X^2=2F^2+2iF\widetilde F$），对 $F$ 配方即得下一式的对偶作用量。原文以“Adding this ... and integrating over $F$”一句带过。
\par
\par\medskip\noindent\hangindent=3.2em\hangafter=1 \textbf{[132]}~\textbf{3.4}~$\tau\to\tau+1$ 对应 $\theta\to\theta+2\pi$：$\theta$ 项为 $\theta(32\pi^2)^{-1}\int F\widetilde F=\theta\nu$，在 $e^{iS}$ 中只依赖 $\theta\bmod 2\pi$，故作用量（路径积分）不变。$T:\tau\to\tau+1$ 与 (3.8) 的 $S:\tau\to-1/\tau$ 一起生成 $SL(2,{\bf Z})$；不同文献对生成元的取法（$S$ 或 $-S$、$T$ 或 $T^{-1}$）只差约定，核对时以下式 (3.11) 的矩阵为准。
\par
\par\medskip\noindent\hangindent=3.2em\hangafter=1 \textbf{[133]}~\textbf{3.4}~$A_D=h$ 与 $h_D=-A$ 的推导：由 $\tau_D=\partial A_D/\partial A=-1/\tau=-\partial A/\partial A_D$，积分即得 $A_D=h(A)$、$h_D(A_D)=-A$（相差可加常数，原文取零）。这正是 $(A_D,A)$ 在 $S:\tau\to-1/\tau$ 下按 (3.11) 的 $S=\begin{pmatrix}0&1\\-1&0\end{pmatrix}$ 变换的含义；代入 $\operatorname{Im}\int d^4\theta\,h\bar A$ 即见其不变（$d^4\theta$ 积分不产生额外因子）。
\par
\par\medskip\noindent\hangindent=3.2em\hangafter=1 \textbf{[134]}~\textbf{3.4}~原文 “equivanlently” 为 “equivalently” 的拼写笔误（原文 PDF 第 52 页）。
\par
\par\medskip\noindent\hangindent=3.2em\hangafter=1 \textbf{[135]}~\textbf{3.4}~此处 $\alpha,\beta,\gamma,\delta$ 是 (3.10) 中 $SL(2,{\rm Z})$ 矩阵的四个元素（即 $\alpha=a,\ \beta=b,\ \gamma=c,\ \delta=d$），与本节其他地方的 $\alpha=g^2/4\pi$ 及 $U(1)_{\cal R}$ 转角无关——同一字母三种含义并存是原文记号的一大陷阱。可自行核对：要求 $\partial{\cal F}'/\partial A'=A'_D$ 就得到 $\beta\gamma=\alpha\delta-1$，即 $\alpha\delta-\beta\gamma=1$。
\par
\par\medskip\noindent\hangindent=3.2em\hangafter=1 \textbf{[136]}~\textbf{3.4}~这段话把度规的 $SL(2,{\rm Z})$ 不变性几何化：$X\cong{\rm C}^2$ 上的辛形式 $\omega={\rm Im}\,da_D\wedge d\bar a$ 在 $SL(2,{\rm Z})$ 下不变（$\operatorname{Im}$ 部分对应 $S$ 保持的反对称形式），而模空间度规是 $\omega$ 沿截面 $f$ 的拉回，故自动不变。“截面”指 $(a_D(u),a(u))$ 作为 $X$ 上丛的局部选取；$u$ 只要求全纯，它与 $\langle{\rm Tr}\phi^2\rangle$ 的最终认同、以及小 $|a|$ 区域的解，将在后续小节确定（原文 PDF 第 52 页）。
\par
\par\medskip\noindent\hangindent=3.2em\hangafter=1 \textbf{[137]}~\textbf{3.4}~“取 $u=a$ 回到原公式”是因为此时 $\frac{da_D}{du}=\tau$，故度规化为 $\operatorname{Im}\tau\,|da|^2$ 即 (3.5)。另注意：$SL(2,{\rm Z})$ 不变性只保证形式，不保证正定性；正是“动能正定”这一物理要求挑出正确的解（见后续小节）。
\par
\par\medskip\noindent\hangindent=3.2em\hangafter=1 \textbf{[138]}~\textbf{3.4}~标题 “Higher Dimensions” 指模空间维数为 $r>1$（更高秩的规范群），不是时空维数；段首的 “\textbackslash{}{\textbackslash{}}” 是作者用来避免 \textbackslash{}verb|\textbackslash{}underline| 后紧跟文字时产生间距的技巧，这里照录。
\par
\par\medskip\noindent\hangindent=3.2em\hangafter=1 \textbf{[139]}~\textbf{3.4}~$f^*(\omega_h)=0$ 等价于 $\partial a_{Di}/\partial a^j=\partial a_{Dj}/\partial a^i$（反对称部分为零），这正是存在局部全纯函数 ${\cal F}(a)$ 使 $a_{Di}=\partial{\cal F}/\partial a^i$ 的积分条件（Poincaré 引理）。原文只说“这一条件保证…”，实际用到的是闭形式局部恰当这一标准事实。
\par
\par\medskip\noindent\hangindent=3.2em\hangafter=1 \textbf{[140]}~\textbf{3.4}~单值变换（monodromy）的定位：把 $(a_D,a)$ 沿模空间中绕奇点的闭合回路作解析延拓，回到起点时一般不再是原来的值，而相差一个 $SL(2,{\bf Z})$ 变换；这正是 $\tau(a)$ 多值性（对数支）的全局体现，后文将用单值矩阵 $M_\infty$ 等拼接各个弱耦合描述。原文 “on being taken around a close loop” 中 “close loop” 应为 “closed loop”（拼写笔误，原文 PDF 第 53 页）。
\par
\par\medskip\noindent\hangindent=3.2em\hangafter=1 \textbf{[141]}~\textbf{3.5}~原文 “microspcopic” 为 “microscopic” 的拼写笔误（原文 PDF 第 53 页）。此处的 $Z$ 是经典（树图）中心荷；下文 (3.12) 给出低能有效理论中的精确形式 $Z=a n_e+a_D n_m$。
\par
\par\medskip\noindent\hangindent=3.2em\hangafter=1 \textbf{[142]}~\textbf{3.5}~这里的质量项可从 $N=2$ 的超势 ${\cal W}=\sqrt{2}\widetilde Q\Phi Q$（见 (3.2) 之前关于超势的讨论）得到：把 $\Phi$ 的标量分量取为真空期望值（原文在此处记作 $A$，通常记作 $a$），即得 $\sqrt{2}n_e a M\widetilde M$，于是 $a\ne0$ 时超多重态获得质量——这是希格斯机制在 $N=2$ 中的实现。
\par
\par\medskip\noindent\hangindent=3.2em\hangafter=1 \textbf{[143]}~\textbf{3.5}~$w=(n_m,n_e)$ 是行矢量，$Z=w^Tv$ 是不变量；等价的列矢量写法是按 $(M^T)^{-1}$ 变换。$M$ 为整系数矩阵是要求电荷格映射到自身，否则 $n_m,n_e$ 不再是整数——这就是“对偶群必须落在 $Sp(2r,{\bf Z})$ 中”的物理含义。
\par
\par\medskip\noindent\hangindent=3.2em\hangafter=1 \textbf{[144]}~\textbf{3.5}~原文此处有一处笔误（原文 PDF 第 54 页）：$\tau_{ij}$ 的变换写作 $(A\tau+B)(c\tau+D)^{-1}$，其中小写 $c$ 与矩阵块 $C$（$M=\begin{pmatrix}A&B\\C&D\end{pmatrix}$）不一致，正确形式应为 $(A\tau+B)(C\tau+D)^{-1}$。同页另有两处小笔误：“geting” 应为 “getting”；$\partial a_Di/\partial a_j$ 的下标应为 $\partial a_{Di}/\partial a_j$。公式与文字已按原文照录，此处如实指出。
\par
\par\medskip\noindent\hangindent=3.2em\hangafter=1 \textbf{[145]}~\textbf{3.6}~原文笔误汇总（本分片）：“associalted”→“associated”（PDF 54）；  “close loop”→“closed loop”（PDF 54、61）；“elementary stat”→“elementary state”（PDF 56）；  “The the associated”→“The associated”（PDF 56）；“gauge theroy”→“gauge theory”、  “in terms our”→“in terms of our”（PDF 58）；“singuralites”→“singularities”（PDF 60）；  “Sine”→“Since”、“anyone”→“any one”（PDF 61）；“regrad”→“regard”、“Upto”→“up to”（PDF 62）；  “in terms complete”→“in terms of complete”（PDF 63）、“hypergeometirc”→“hypergeometric”（PDF 64）。  译文一律按正确拼写处理；所有公式照原文抄录，未作“修正”。
\par
\par\medskip\noindent\hangindent=3.2em\hangafter=1 \textbf{[146]}~\textbf{3.6}~式 (\ref{mondone})（原文编号 (97)，原文 PDF 第 54 页）有三处易错：  (1) 单值矩阵作用在列向量 $(a_D,a)^T$ 上，而电荷向量按 $(M^{-1})^T$ 变换（见原文 PDF 第 54 页      式 (97) 之前关于 $w=(n_m,n_e)\to wM^{-1}$ 的说明）；两者方向相反，写反会得到 $M_\infty^{-1}$。  (2) $M_\infty=PT^{-2}$ 中 $P=-I$ 是 $SL(2,{\rm Z})$ 的中心元，给出整体变号      $a\to -a,\ a_D\to -a_D$；$T^{-2}$ 给出 $a_D$ 的平移 $2a$。二者不可拆错。  (3) 这里的回路是绕 $u=0$ 的大回路，它与绕 $u=\infty$ 的回路同伦（$u$ 平面上两者只差一块可缩区域），      故记为 $M_\infty$；绕行方向取逆时针，若取顺时针则所有单值矩阵都取逆。
\par
\par\medskip\noindent\hangindent=3.2em\hangafter=1 \textbf{[147]}~\textbf{3.6}~原文 PDF 第 55 页第一段。$M_\infty$ 把 $a\to -a$，故 $a^2$ 在 $M_\infty$ 下不变；  若其余单值变换都与 $M_\infty$ 对易，单值群便成为阿贝尔群，$a^2$ 才可能是整体坐标。但此时  $\tau=da_D/da$ 的虚部无法保持正定，动能项会变号——“非阿贝尔的单值群”与“度规正定”是同一件事  的两面；这正是原文断言“this cannot be the case”的物理内容，原文未展开。
\par
\par\medskip\noindent\hangindent=3.2em\hangafter=1 \textbf{[148]}~\textbf{3.6}~原文 PDF 第 55 页。严格说，威尔逊有效作用量是把有质量自由度积掉得到的；当某个质量  趋于零时，该自由度便不应再被积掉，有效作用量的解析结构在此失效，表现为模空间上的奇点。奇点的  单值变换由在奇点处变为无质量的粒子的量子数决定——这正是后文用 $(n_m,n_e)$ 标记 $M(n_m,n_e)$ 的原因。  术语 “integrate out” 译作“积掉”。
\par
\par\medskip\noindent\hangindent=3.2em\hangafter=1 \textbf{[149]}~\textbf{3.6}~式 (\ref{mondq})（原文编号 (98)，原文 PDF 第 56 页）。与式 (97) 对比：在大 $u$  （弱耦合）处 $a_D\to -a_D+2a$、$a\to -a$；在超多重态无质量点处 $a_D\to a_D+2a$、$a\to a$。  两者在 PSL(2,Z) 意义下只差一个整体 $P=-I$，即 $M_\infty$ 对应 $T^{-2}$。需注意 $-I\in\Gamma(2)$  （因为 $-1\equiv 1\pmod 2$），所以把 $M_\infty$ 写成 $PT^{-2}$ 并不影响“单值群由 $M_1,M_{-1}$  生成并等于 $\Gamma(2)$”这一结论；但若在模变换中忽略这个 $-I$，就会误把 $(a_D,a)$ 与  $(-a_D,-a)$ 当成同一个截面。
\par
\par\medskip\noindent\hangindent=3.2em\hangafter=1 \textbf{[150]}~\textbf{3.6}~式 (\ref{asns})（原文编号 (99)，原文 PDF 第 56 页）。$SL(2,{\rm Z})$ 矩阵  $M=\left(\begin{smallmatrix}\alpha&\beta\\ \gamma&\delta\end{smallmatrix}\right)$ 作用在周期向量上为  $v\to Mv$；由 $Z=n_ma_D+n_ea$ 不变可知电荷向量按 $(M^{-1})^T$ 变换，即  $(n_m,n_e)\to(\delta n_m-\gamma n_e,\,-\beta n_m+\alpha n_e)$。原文第一分量写作 $n_m\delta-n_e\gamma$，  核对时极易把 $\gamma$ 与 $n_e$ 的配对弄错；第二分量同理。另注意 $\alpha\delta-\beta\gamma=1$  正是 $SL(2,{\rm Z})$ 的行列式条件。
\par
\par\medskip\noindent\hangindent=3.2em\hangafter=1 \textbf{[151]}~\textbf{3.6}~式 (\ref{mondn})（原文编号 (100)，原文 PDF 第 56 页）。${\rm Tr}\,M(n_m,n_e)=2$ 且  $M\ne\pm I$，故 $M$ 是 $SL(2,{\rm Z})$ 中的抛物型元素，必与 $T$ 共轭；这是“只需两个奇点就能生成  整个单值群”的群论根据。矩阵元 $2n_e^2$、$-2n_m^2$ 中的因子 2 来自单圈 $\beta$ 函数（见下一小节），  不是笔误。
\par
\par\medskip\noindent\hangindent=3.2em\hangafter=1 \textbf{[152]}~\textbf{3.6}~原文 PDF 第 56 页式 (101) 之后的解集。由矩阵等式得到四个方程，其中 $m^2=m'^2$  只给出 $m=\pm m'$，还需与 $1+mn=m'n'+2m'^2$、$1-mn=-m'n'$ 联立才能定出 $m,m'=\pm 1$。  表中四列依次对应 $(m,m')=(1,1),(1,-1),(-1,-1),(-1,1)$，最简单的解 $m=m'=1,n=0,n'=-1$  是第一列。另注意矩阵等式的左端是 $M_1$、右端是 $M_\infty M_{-1}^{-1}$，与式 (101) 的乘积次序  $M_1M_{-1}=M_\infty$ 一致；若把次序写反，所得方程会变成另一组。
\par
\par\medskip\noindent\hangindent=3.2em\hangafter=1 \textbf{[153]}~\textbf{3.6}~式 (\ref{mon-inv})（原文编号 (102)，原文 PDF 第 57 页）之后。原文 “Clearly,  $q_m=m, q_n=n$ is a solution” 中 $q_n$ 应为 $q_e$（本征向量按 $(q_m,q_e)$ 排列，见式 (102) 左端）；  这是原文笔误，公式本身无误。该本征值方程只定出 $(q_m,q_e)$ 的方向：$nq_m-mq_e=0$ 的通解为  $(q_m,q_e)\propto(m,n)$；对稳定双子取最小整数倍，才得到唯一的 $(m,n)$。
\par
\par\medskip\noindent\hangindent=3.2em\hangafter=1 \textbf{[154]}~\textbf{3.6}~原文 PDF 第 57 页末。$M_\infty$ 对电荷的作用是  $(n_m,n_e)\to(-n_m,-n_e-2n_m)$（见第 54 页），故在磁荷固定的扇区里，电荷只在模 $2n_m$  的意义下定义；对 $n_m=1$ 就是模 2。用 $M_\infty^n$ 共轭只改变基本群表示的基底（即单值矩阵的  起点），不改变物理；但它会改变“哪个双子表现为基本态”的标记，后文取 $M_1,M_{-1}$ 的表示时  已固定了这一基底。
\par
\par\medskip\noindent\hangindent=3.2em\hangafter=1 \textbf{[155]}~\textbf{3.7}~原文 PDF 第 57--58 页 $\beta$ 函数一节。这里 $b$ 定义为 $\beta(g)=b\,g^3$ 中 $g^3$  的系数（已含 $1/16\pi^2$），故 $\mu\,d(1/\alpha)/d\mu=-8\pi b$ 中的因子是 $8\pi$ 而不是 $2\pi$；  由 $\alpha=g^2/4\pi$ 可直接验证：$d(1/\alpha)/d\ln\mu=-(8\pi/g^3)\beta(g)$。系数 $2/3$ 与 $1/3$  分别是外尔费米子与复标量的标准单圈贡献，最易抄错。
\par
\par\medskip\noindent\hangindent=3.2em\hangafter=1 \textbf{[156]}~\textbf{3.7}~原文 PDF 第 58 页。一个超多重态含两个外尔费米子、两个复标量，故  $2\cdot\frac{2}{3}+2\cdot\frac{1}{3}=2$，得 $b=Q^2/8\pi^2$。这里每个外尔费米子贡献  $2Q_f^2/3$、每个复标量贡献 $Q_s^2/3$（均以 $g^3/16\pi^2$ 为单位）；$b>0$ 表明该 $U(1)$  在紫外并不渐近自由，而是红外自由（紫外出现朗道极点）。
\par
\par\medskip\noindent\hangindent=3.2em\hangafter=1 \textbf{[157]}~\textbf{3.7}~原文 PDF 第 58 页。用费米子的手征转动把 $\theta$ 参数转掉，前提是只关心 $\theta=0$  的真空；一般情形 $\tau=\theta/2\pi+i/\alpha$，$\mu\,d\tau/d\mu=-iQ^2/\pi$ 只对虚部成立。此处  “反常”指 $U(1)_R$ 的手征反常，它使 $\theta$ 的移动可完全转移到费米子相位上。对磁单极多重态，  把上述计算整体搬到对偶变量即得 $\tau_D\simeq-(i/\pi)\ln(a_D/\Lambda)$；注意 $\tau_D$ 与 $\tau$  不是同一个函数，在强耦合区不能混用。
\par
\par\medskip\noindent\hangindent=3.2em\hangafter=1 \textbf{[158]}~\textbf{3.8}~原文 PDF 第 58 页（本小节开头）有两处笔误：“gauge theroy” 应为 “gauge theory”，  “in terms our” 应为 “in terms of our”。另注意：加 $W=m\,{\rm Tr}\Phi^2$ 使 $N=2$ 破缺到 $N=1$；  低能纯阿贝尔理论被认为有质量间隙，并自发破缺 $\mathbf{Z}_4\to\mathbf{Z}_2$，这是 $N=1$  规范理论禁闭的宏观图像，下文用磁单极凝聚来“实现”它。
\par
\par\medskip\noindent\hangindent=3.2em\hangafter=1 \textbf{[159]}~\textbf{3.8}~原文 PDF 第 59 页。由 $\hat W=\sqrt2 A_DM\widetilde M+mU(A_D)$ 得  $\partial\hat W/\partial M=\sqrt2 A_D\widetilde M$、  $\partial\hat W/\partial\widetilde M=\sqrt2 A_DM$、  $\partial\hat W/\partial A_D=\sqrt2 M\widetilde M+m\,du/dA_D$，三式令零即得文中两行方程  （第二行把两个方程合并书写，原文用小写 $a_D$ 指其标量分量）。$M=\widetilde M=(-m u'(0)/\sqrt2)^{1/2}$  中的 $u'(0)$ 指 $du/da_D$ 在 $a_D=0$ 处的值；原文未讨论根号下为负（此时取纯虚相位）的约定。这一相位只影响凝聚体的相位，$M\ne0$（从而磁希格斯机制成立）的结论不变。
\par
\par\medskip\noindent\hangindent=3.2em\hangafter=1 \textbf{[160]}~\textbf{3.9}~图 3（原文 PDF 第 60 页）与原文 PDF 第 59--60 页的丛语言。$u$ 平面去掉三个奇点后，  周期向量 $(a_D,a)^T$ 是平坦 $SL(2,{\rm Z})$ 丛 $V$ 的全纯截面；绕 $u=-1$、$u=1$、$u=\infty$ 的  回路分别给出 $M_{-1},M_1,M_\infty$，且按回路的复合次序 $M_\infty=M_1M_{-1}$（先绕 $u=-1$，  再绕 $u=1$，合成回路等效于绕无穷远）。把乘积次序写反会得到 $M_\infty^{-1}$。
\par
\par\medskip\noindent\hangindent=3.2em\hangafter=1 \textbf{[161]}~\textbf{3.9}~原文 PDF 第 60 页未编号的渐近行为公式。弱耦合下 $u=\frac12a^2$（原文第 54 页），  故 $a\cong\sqrt{2u}$；而 $a_D\approx i\frac{\sqrt{2u}}{\pi}\ln u$ 中的 $i$ 与 $\pi$ 来自单圈预势  （见第 54 页 $\partial{\cal F}/\partial a$ 一式）。注意此式是对 $u$ 取对数，而不是对 $a$ 取对数，  二者相差 $\ln a^2=2\ln a$ 的常数项。在 $u=1$ 附近，$a_D\approx c_0(u-1)$ 是“磁单极无质量”的标志，  而 $a\approx a_0+\frac{i}{\pi}a_D\ln a_D$ 表明此时 $a_D$ 才是好坐标、$a$ 是坏坐标；$u=-1$ 处  作替换 $a_D\to a-a_D$。
\par
\par\medskip\noindent\hangindent=3.2em\hangafter=1 \textbf{[162]}~\textbf{3.9}~式 (\ref{tau-metric})（原文编号 (105)，原文 PDF 第 60 页）。$\tau=da_D/da$ 在 $u$  坐标下写成 $(da_D/du)/(da/du)$；弱耦合区 $\tau=\theta/2\pi+4\pi i/g^2$，故  ${\rm Im}\,\tau=4\pi/g^2>0$ 正是动能项正定。$S$ 对偶 $\tau\to-1/\tau$ 伴随  $(a_D,a)\to(a,-a_D)$ 的 $SL(2,{\rm Z})$ 作用，即“电—磁”互换；不要把 $a_D$ 与 $a$ 误当成纯虚、  纯实的量，它们一般是复的，只有 $\tau$ 的虚部才有正定意义（原文此处对“虚实变换”未加说明）。
\par
\par\medskip\noindent\hangindent=3.2em\hangafter=1 \textbf{[163]}~\textbf{3.9}~式 (\ref{cubic})（原文编号 (106)，原文 PDF 第 60 页）。验证 $\mathbf{Z}_4$：在  $u\to-u,\ x\to-x,\ y\to\pm iy$ 下，右端变为 $(-x-1)(-x+1)(-x+u)=-(x^2-1)(x-u)$，而左端  $(\pm iy)^2=-y^2$，故方程不变。其中 $u\to-u$ 与 $x\to-x$、$y\to\pm iy$ 合起来才构成  $\mathbf{Z}_4$，而只有 $u\to-u$ 这一 $\mathbf{Z}_2$ 子群作用在模空间上；曲线族 $E_u$ 本身  则具有更大的 $\mathbf{Z}_4$。
\par
\par\medskip\noindent\hangindent=3.2em\hangafter=1 \textbf{[164]}~\textbf{3.9}~图 4、图 5（原文 PDF 第 61 页）。$y^2=(x-1)(x+1)(x-u)$ 对 $y$ 是二次的，故  $x$ 空间是 $\mathbf{C}\cup\{\infty\}$ 的双重覆叠，四个分支点 $-1,1,u,\infty$ 两两相连成两条割线；  割线的具体取法不唯一（这里取 $[-1,1]$ 与 $[u,\infty]$），不同取法只改变同调基的辛基底，不改变  物理。图 5 用环面说明：$c_1,c_2$ 压扁为 $l_1,l_2$ 后剪开，得到双叶结构；$x$ 平面上绕单条割线的  回路对应 $a$-圈，与两条割线都相交的回路对应 $b$-圈。
\par
\par\medskip\noindent\hangindent=3.2em\hangafter=1 \textbf{[165]}~\textbf{3.9}~原文 PDF 第 62 页。奇点处消失圈与单值变换的对应是本节的核心：$u=\infty$ 时  $a$-圈收缩（对应 $M_\infty$），$u=1$ 时 $b$-圈收缩（对应 $M_1$）；由 $\mathbf{Z}_2$ 对称性  $u\to-u$，$u=-1$ 处同样是一个 $b$-圈收缩（对应 $M_{-1}$，但基底经 $a_D\to a-a_D$ 重排）。  这也解释了 $H/\Gamma(2)$ 为何恰有三个尖点——尖点即消失圈的位置。
\par
\par\medskip\noindent\hangindent=3.2em\hangafter=1 \textbf{[166]}~\textbf{3.9}~原文 PDF 第 62 页。原文写 “first homology group $V_u=H^1(E_u,{\rm C})$”，严格说  $H^1$ 是上同调记号（同调应为 $H_1$）；因 $E_u$ 是紧致定向曲面，二者经 Poincaré 对偶同构，  后文按同调使用无碍，但引用时宜留意这一记号混用。留数为零的条件保证 $\oint_\gamma\lambda$  在 $\gamma$ 连续形变时不变（即使越过 $\lambda$ 的极点），这是 $\lambda$ 能作为 $V_u$ 元素并  与周期配对的前提。
\par
\par\medskip\noindent\hangindent=3.2em\hangafter=1 \textbf{[167]}~\textbf{3.9}~原文 PDF 第 63 页。由 $d\lambda/du=f(u)\,dx/y$ 与 $f(u)=-\sqrt2/4\pi$ 积分可得  $\lambda=\frac{\sqrt2}{2\pi}\frac{dx\sqrt{x-u}}{\sqrt{x^2-1}}$：因为  $\int du\,(x-u)^{-1/2}=-2\sqrt{x-u}$，负号与 $f$ 的负号相消，系数恰为 $\sqrt2/2\pi$。三个等号  用到 $y^2=(x^2-1)(x-u)$，即 $\frac{\sqrt{x-u}}{\sqrt{x^2-1}}=\frac{y}{x^2-1}=\frac{x-u}{y}$。  $\sqrt{x-u}$ 的分支约定与下面 $a_D$ 积分中出现的 $i$ 直接相关，不能随意选取。
\par
\par\medskip\noindent\hangindent=3.2em\hangafter=1 \textbf{[168]}~\textbf{3.9}~式 (\ref{au})、(\ref{adu})（原文编号 (107)、(108)，原文 PDF 第 63 页）。$\gamma_2$  沿 $[-1,1]$ 割线来回，故 $a(u)$ 的积分限为 $-1\to 1$；$\gamma_1$ 绕 $1$ 与 $u$，故积分限为  $1\to u$。两式的积分路径都取在相应割线的同一侧，上下限方向决定整体符号；若把 $\gamma_1$ 的  定向取反，$a_D$ 变号，继而 $M_1$、$M_{-1}$ 都会改变，务必与原文约定一致。原文在此引 \cite{SW-I}  验证了渐近行为。
\par
\par\medskip\noindent\hangindent=3.2em\hangafter=1 \textbf{[169]}~\textbf{3.9}~原文 PDF 第 64 页超几何函数的两种表示。积分表示前的 $\Gamma$ 因子是  $\Gamma(\gamma)/[\Gamma(\beta)\Gamma(\gamma-\beta)]$，而级数表示前是  $\Gamma(\gamma)/[\Gamma(\alpha)\Gamma(\beta)]$——两者并不相同，极易误抄成同一个；积分表示还要求  ${\rm Re}\,\gamma>{\rm Re}\,\beta>0$。另注意原文把 $F$ 的自变量写在分号之后  （$F(\alpha,\beta,\gamma;z)$），与某些文献的 $F(\alpha,\beta;\gamma;z)$ 记法只是标点差别。
\par
\par\medskip\noindent\hangindent=3.2em\hangafter=1 \textbf{[170]}~\textbf{3.9}~原文 PDF 第 64 页。在 $a_D$ 的积分 (\ref{adu}) 中作 $x=(u-1)t+1$，则  $x-u=(u-1)(t-1)$；在 $u>1$ 的割线一侧、$t\in[0,1]$ 上，  $\sqrt{x-u}=i\sqrt{u-1}\sqrt{1-t}$，同时  $\sqrt{x^2-1}=\sqrt{(u-1)t\,[(u-1)t+2]}$。整理后出现的 $i$、$(u-1)$ 以及  $(1-u)/2=-k'^2/k^2$ 中的负号，是这一段最易出错的地方，务必逐步核对。
\par
\par\medskip\noindent\hangindent=3.2em\hangafter=1 \textbf{[171]}~\textbf{3.9}~原文 PDF 第 64 页椭圆积分约定。这里 $K(k),E(k)$ 的宗量是模数 $k$，而不是参数  $m=k^2$：$K(k)=\frac{\pi}{2}F(1/2,1/2,1;k^2)$，$E(k)=\frac{\pi}{2}F(-1/2,1/2,1;k^2)$。原文只写  $E(k),K(k)$ 而未声明约定，引用文献数据时务必注意 $K(m)$ 与 $K(k)$ 相差一个 $k\to\sqrt m$  的代换；$k'^2=1-k^2$ 为补模。
\par
\par\medskip\noindent\hangindent=3.2em\hangafter=1 \textbf{[172]}~\textbf{3.9}~原文 PDF 第 65 页。$F(a,b,c;z)=(1-z)^{-a}F(a,c-b,c;z/(z-1))$ 是超几何函数的分式  线性变换，由此得 $K(ik'/k)=kK'(k)$、$E(ik'/k)=k^{-1}E'(k)$。注意 $|ik'/k|$ 一般不在 $[0,1)$  内，$K,E$ 需解析延拓；又 $E'$、$K'$ 是 $E(k')$、$K(k')$ 的简写，不是对 $k$ 的导数——原文的  撇号在这两节里有“补模”与“导数”两种含义，读到 $dE'/dk$ 时尤须分清。
\par
\par\medskip\noindent\hangindent=3.2em\hangafter=1 \textbf{[173]}~\textbf{3.9}~原文 PDF 第 65 页末。由 $a=\frac{4}{\pi k}E(k)$、  $a_D=\frac{4}{\pi i}\frac{E'-K'}{k}$ 对 $k$ 求导：  $\frac{da}{dk}=-\frac{4}{\pi}\frac{K}{k^2}$，  $\frac{da_D}{dk}=\frac{4}{\pi i}\frac{K'}{k^2}$（注意 $E'=E(k')$、$K'=K(k')$ 中 $k'^2=1-k^2$  也依赖 $k$，求导时不能漏掉 $dk'/dk=-k/k'$），二者相除才得 $\tau=iK'/K$。若漏掉补模对 $k$ 的  依赖，会得到相差一个 $k$ 因子的结果。
\par
\par\medskip\noindent\hangindent=3.2em\hangafter=1 \textbf{[174]}~\textbf{4}~见原文 PDF 第 65 页导语第一段。原文称“为保证渐近自由要求 $N_f\le4$”， 但由式 (111) 可见 $N_f=4$ 时 $\beta(g)=0$：该理论是有限的（共形不变）， 并非渐近自由；严格的渐近自由条件是 $N_f<4$。此处是原文的宽泛表述，阅读时不要 把 $N_f=4$ 当作渐近自由理论。
\par
\par\medskip\noindent\hangindent=3.2em\hangafter=1 \textbf{[175]}~\textbf{4}~见原文 PDF 第 65 页。与第 3 节的纯规范归一化对照：那里 $a\cong\sqrt{2u}$、基本表示的夸克 $n_e=\pm1/2$、$W$ 玻色子 $n_e=\pm1$； 本节把 $n_e$ 乘 2、$a$ 除 2，于是 $a\cong\frac12\sqrt{2u}$ 而 $a_D$ 不变 （$2a_D=\partial{\cal F}/\partial a$）。这个因子 2 贯穿全节，尤其是式 (113)、第 76 页的换基与式 (115)，核对时最易出错。
\par
\par\medskip\noindent\hangindent=3.2em\hangafter=1 \textbf{[176]}~\textbf{4.1}~见原文 PDF 第 66 页。$A$ 是光子矢量多重态的 $N=1$ 手征超场 （其最低分量为复标量），$W_\alpha$ 是场强手征超场；$M^i$ 与 $\widetilde M_i$ 的味指标 $i$ 与 $U(1)$ 规范指标都隐含缩并。超势中的 $\sqrt2$ 系数不是可以 随意改动的约定：它保证同一超多重态内玻色子与费米子质量简并，并与 BPS 质量 公式中的 $\sqrt2$ 相关联。
\par
\par\medskip\noindent\hangindent=3.2em\hangafter=1 \textbf{[177]}~\textbf{4.1}~见原文 PDF 第 66 页。原文 “afflicted by an anomaly” 一语指该整体 对称性被瞬子反常破缺（下一小节将具体算出 ${\bf Z}_{4N_c-2N_f}$），并非说它 完全不是对称性；故译作“受反常的破坏”，而非“具有反常”。
\par
\par\medskip\noindent\hangindent=3.2em\hangafter=1 \textbf{[178]}~\textbf{4.1}~见原文 PDF 第 67 页式 (109)。QED 的 $U(1)$ $\beta$ 函数为正（红外自由），故 $a_D$ 的系数是 $-ik/2\pi$，与纯 $N=2$ 理论大 $u$ 处的 $+i(4-N_f)/2\pi$（式 (113)）反号；此处极易 漏掉负号。另注意该式是对 $k$ 个电子求和前的单圈结果，$\ln(a/\Lambda)$ 的 宗量是 $a$（量纲 1），与式 (113) 中的 $u/\Lambda^2_{N_f}$ 不同。 由 $\tau=\partial a_D/\partial a$ 可核验 ${\rm Im}\,\tau=0$ 恰好发生在 $|a|=\Lambda/e$，即朗道极点。
\par
\par\medskip\noindent\hangindent=3.2em\hangafter=1 \textbf{[179]}~\textbf{4.1}~见原文 PDF 第 67 页末。中心荷多出 $\sum_iS_im_i/\sqrt2$ 项， 故 BPS 质量 $M=\sqrt2|Z|$ 给出 $|\sqrt2a+m_i|$，而不是 $\sqrt2|a|$； 核对时注意 $m_i$ 前的 $1/\sqrt2$ 与 $a$ 前的 $\sqrt2$ 来自同一归一化， 不可对调。后文式 (84)（原书第 84 式）即此修改后的质量公式。
\par
\par\medskip\noindent\hangindent=3.2em\hangafter=1 \textbf{[180]}~\textbf{4.2}~见原文 PDF 第 67--68 页。$SU(2)$ 中基本表示与共轭表示同构， $Q^i$ 与 $\widetilde Q_i$ 才能拼成 $O(2N_f)$ 的 $2N_f$ 维矢量；对 $SU(N_c)$（$N_c>2$）这一步不成立，故 $N_c=2$ 的味对称性比一般情形大。 关于商群：$U(1)_{\cal R}$ 中心的 $Z_2$、$SU(2)_R$ 中心的 ${\bf Z}_2$ 与 $O(2N_f)$ 中心的 ${\bf Z}_2$ 是同一个 ${\bf Z}_2$ 的不同化身，故整体 对称群要按对角线 ${\bf Z}_2$ 取商，否则会重复计数。
\par
\par\medskip\noindent\hangindent=3.2em\hangafter=1 \textbf{[181]}~\textbf{4.2}~见原文 PDF 第 68 页式 (111)。负号对应 $N_f<4$ 的 渐近自由；$N_f=4$ 时 $\beta=0$，理论有限而非渐近自由（见本节导语的批注）。 若改用 $\alpha=g^2/4\pi$，渐近自由条件同样是 $4-N_f>0$。系数 $1/16\pi^2$ 与 $g$ 的归一化有关，比较不同文献时先核对其 $\beta$ 函数定义。
\par
\par\medskip\noindent\hangindent=3.2em\hangafter=1 \textbf{[182]}~\textbf{4.2}~见原文 PDF 第 68--69 页。$N_c=2$ 时 ${\bf Z}_{2(4N_c-2N_f)}={\bf Z}_{2(8-2N_f)}={\bf Z}_{16-4N_f} ={\bf Z}_{4(4-N_f)}$，指数中 $4N_c-2N_f=8-2N_f$ 与后文 $4(4-N_f)$ 只有在 $N_c=2$ 时才一致，代值核对时勿混用。$N_f=0$ 时没有夸克抵消反常， 只有变换的平方无反常，这与第 3 节的讨论一致。
\par
\par\medskip\noindent\hangindent=3.2em\hangafter=1 \textbf{[183]}~\textbf{4.2}~见原文 PDF 第 69 页。$\omega={\exp}(2i\pi/4(4-N_f))$ 应理解为 $\omega=e^{2\pi i/[4(4-N_f)]}$，即 ${\bf Z}_{4(4-N_f)}$ 的本原根；原文排版 “$2i\pi/4(4-N_f)$” 极易被误读成 $(2\pi i/4)(4-N_f)$。可用下一段的结果反验： $u$ 的荷为 4，故 $\omega^4=e^{2\pi i/(4-N_f)}$，恰给出 $u\to e^{2\pi i/(4-N_f)}u$。
\par
\par\medskip\noindent\hangindent=3.2em\hangafter=1 \textbf{[184]}~\textbf{4.2}~原文 PDF 第 69 页式 (112) 第二组第二行左端写作 $\widetilde{Q}_1$，据上下文（该组对 $i\ne1$ 的场统一乘 $\omega^{-1}$， 且与第一组的 $Q^1\leftrightarrow\widetilde Q_1$ 置换相区别）应为 $\widetilde{Q}_i$。这是原文的排印笔误，公式按原样保留，阅读时请自行修正。
\par
\par\medskip\noindent\hangindent=3.2em\hangafter=1 \textbf{[185]}~\textbf{4.2}~见原文 PDF 第 69 页。$a$ 的展开从 $n=1$ 开始，而 $a_D$ 的 展开从 $n=0$ 开始（$n=0$ 项是 $\sqrt u\,a_{D0}(N_f)$，即解析的非微扰修正）； 两式的幂次都是 $(\Lambda^2_{N_f}/u)^{n(4-N_f)}$，其中 $\Lambda^2_{N_f}$ 与 $u$ 同量纲。注意 $a_D$ 对数项系数 $i(4-N_f)/(2\pi)$ 与 $a$ 的领头项 $\frac12\sqrt{2u}$ 的 $N_f$ 依赖不同，且因子 $\sqrt u$ 整体提到求和号外， 抄写时最易漏掉。
\par
\par\medskip\noindent\hangindent=3.2em\hangafter=1 \textbf{[186]}~\textbf{4.2}~见原文 PDF 第 69 页。把 $\Lambda_{N_f}^{4-N_f}$ 的 $U(1)_{\cal R}$ 荷定为 $2(4-N_f)$、$\rho$ 宇称为奇，这一赋值在 $N_f=4$ 时 退化为 $\Lambda^0$ 而失去意义；所以原文随后把 $N_f=4$ 单独处理。也因此， $N_f=4$ 的曲线不能由上述 $a,a_D$ 展开式直接取极限得到（参见原文第 79--80 页 与 [32]）。
\par
\par\medskip\noindent\hangindent=3.2em\hangafter=1 \textbf{[187]}~\textbf{4.3}~见原文 PDF 第 70 页。要点是“$N_f$ 个超多重态 $\Rightarrow$ $2N_f$ 个费米子二重态 $\Rightarrow$ 每个二重态在磁单极背景中各贡献一个 零模 $\Rightarrow$ 共 $2N_f$ 个零模 $\Rightarrow$ $2N_f$ 维 Clifford/Dirac 代数 $\Rightarrow$ $Spin(2N_f)$ 的旋量”。原文用“each $SU(2)$ doublet of fermions has one zero-mode”一句带过零模个数的计算（指标定理的结论）， 具体可参看第 2、3 节关于零模的讨论。
\par
\par\medskip\noindent\hangindent=3.2em\hangafter=1 \textbf{[188]}~\textbf{4.3}~见原文 PDF 第 70 页。$e^{i\pi Q}=e^{i\theta n_m}(-1)^H$ 中 “$2\pi$ 转动不是恒等变换”来自 Witten 效应：绕磁单极一周的电荷转动同时 产生一个 $\theta$ 角的规范变换。$(-1)^H$ 就是 $SU(2)$ 的中心元， 等价于 $Spin(2N_f)$ 旋量表示的手征算符，故 $n_e$ 的奇偶性与 $SO(2N_f)$ 手征性相关联。核对时注意 $n_e$ 的符号约定（本节已把 $n_e$ 乘 2， 基本夸克 $n_e=\pm1$，见第 65 页）。
\par
\par\medskip\noindent\hangindent=3.2em\hangafter=1 \textbf{[189]}~\textbf{4.3}~见原文 PDF 第 71 页。三条记号极易混淆：$N_f=2$ 时中心为 ${\bf Z}_2\times{\bf Z}_2$，夸克的荷为 $(1,1)$，磁单极贡献使第一个 ${\bf Z}_2$ 指标变为 $(n_e+n_m)\bmod2$；$N_f=3$ 时 $Spin(6)=SU(4)$ 的 中心是 ${\bf Z}_4$，基本夸克在 $6$ 中、相对中心的荷为 2，故相位是 $e^{i\pi(n_m+2n_e)/4}$；$N_f=4$ 时 $Spin(8)$ 的中心两个 ${\bf Z}_2$ 分别给出 $n_m$、$n_e$ 的奇偶。原文在 $N_f=4$ 处写“representatios” （应为 representations），系排印笔误。
\par
\par\medskip\noindent\hangindent=3.2em\hangafter=1 \textbf{[190]}~\textbf{4.3}~见原文 PDF 第 71 页。这里 $S_{N_f}$ 是 $SO(2N_f)$ 破缺后 剩余 $SO(2)$ 对应的整体荷，$m_{N_f}/\sqrt2$ 是它耦合到中心荷的方式； 由 $M=\sqrt2|Z|$ 得夸克质量 $|\sqrt2a+m_{N_f}|$，故在 $a=\pm m_{N_f}/\sqrt2$ 处变无质量。符号约定沿用式 (84) 的 修改后质量公式，比较时注意 $S_{N_f}$ 的归一化。
\par
\par\medskip\noindent\hangindent=3.2em\hangafter=1 \textbf{[191]}~\textbf{4.4}~见原文 PDF 第 72 页。印出的 ${\cal M}^{-1}$ 第三行是 $(0,1,1)$，这是原文笔误；按 ${\cal M}$ 直接求逆应为 $\left(\begin{array}{ccc}1&0&0\\1&1&-1\\0&0&1\end{array}\right)$。 用紧接着的电荷变换律反验即知：$W=(S,n_m,n_e)$、$W\to W{\cal M}^{-1}$ 若按 印出的矩阵计算会得到 $(S+n_m,\;n_m+n_e,\;-n_m+n_e)$，其中第二个分量多出 $+n_e$，与下一段给出的 $(S+n_m,n_m,-n_m+n_e)$ 矛盾。公式按原文保留， 阅读时请自行修正。
\par
\par\medskip\noindent\hangindent=3.2em\hangafter=1 \textbf{[192]}~\textbf{4.4}~见原文 PDF 第 72 页。此逆矩阵表达式分子上没有显式的 $1/\det{\cal M}$，是因为此处 $\det{\cal M}=1$（原文随后即声明）。若把 $\det{\cal M}=1$ 误当成 $\det B=1$（$B$ 为右下 $2\times2$ 块），或忘记 下三角块的符号，都会得到错误结果；推导时写成 ${\cal M}=\left(\begin{smallmatrix}1&0\\ b&B\end{smallmatrix}\right)$、 ${\cal M}^{-1}=\left(\begin{smallmatrix}1&0\\ -B^{-1}b&B^{-1}\end{smallmatrix}\right)$ 最稳妥。
\par
\par\medskip\noindent\hangindent=3.2em\hangafter=1 \textbf{[193]}~\textbf{4.4}~见原文 PDF 第 72--73 页。这一段是“边缘稳定性”（marginal stability）论证的雏形：当两个态的中心荷相位相同时，单粒子态与多粒子态 可以简并并相互衰变，因此绕奇点一周后单粒子态可能变成多粒子态。原文说 “the formalism is $SL(2,Z)$ covariant, the spectrum is not”，要点是 单值变换只保证 BPS 质量公式（中心荷）协变，不保证谱的粒子数不变； 该论证在 Seiberg--Witten 原文 [32] 中有更详细的讨论。
\par
\par\medskip\noindent\hangindent=3.2em\hangafter=1 \textbf{[194]}~\textbf{4.4}~见原文 PDF 第 73 页。$N_f=4$ 时对偶群不是单纯的 $SL(2,Z)$， 而是半直积 $Spin(8)\rtimes SL(2,Z)$：$SL(2,Z)$ 单独作用会把矢量表示 $v$ 与旋量表示 $s,c$ 互换（因为 $SL(2,Z)$ 变换改变 $(n_m,n_e)$ 的奇偶性）， 必须同时用三元性把 $v,s,c$ 再置换回来。原文“$SL(2,Z)$ alone cannot keep this spectrum unchanged”即指此。类似地，$N_f<4$ 的有物质理论中， 单值变换含 $S$ 荷的平移，故对偶群一般也不是 $SL(2,Z)$（见式 (115) 前的 $3\times3$ 矩阵讨论）。
\par
\par\medskip\noindent\hangindent=3.2em\hangafter=1 \textbf{[195]}~\textbf{4.5}~见原文 PDF 第 73 页。匹配条件的指数极易记错：$\Lambda_{N_f}$ 在 $\Lambda_{N_f}^{4-N_f}$ 中出现，故对 $N_f=3$ 是 $\Lambda_3$ 的一次方 （$\Lambda_0^4=m^3\Lambda_3$），而对 $N_f=1$ 是 $\Lambda_1^3$ （$\Lambda_0^4=m\Lambda_1^3$）。注意这里的 $\Lambda_{N_f}$ 是单重标度 （量纲 1），与式 (113) 中出现的 $\Lambda^2_{N_f}$ 不是同一 写法，核对时先统一记号。
\par
\par\medskip\noindent\hangindent=3.2em\hangafter=1 \textbf{[196]}~\textbf{4.5}~见原文 PDF 第 74 页。$N_f=3$ 的计数：大 $m$ 时有三个奇点 （荷为 $\underline{3},\underline{1},\underline{1}$），$m\to0$ 时前两个 合并成 $SU(4)$ 的 $\underline{4}$（$\underline{3}+\underline{1}$ 恰好 给出 $\underline{4}$），第三个保持为 $\underline{1}$。$\underline{1}$ 对应的 $(n_m,n_e)=(2,1)$ 是磁单极--磁单极束缚态，原文明确说半经典图像 只保证第一个态存在，第二个的存在性未定；这是原文承认的开放性假设， 不应在译文中当作已证结论。
\par
\par\medskip\noindent\hangindent=3.2em\hangafter=1 \textbf{[197]}~\textbf{4.5}~见原文 PDF 第 74 页。注意与上一小节式 (110) 之后 的 $(\epsilon,\epsilon')$ 公式对照：那里第一个指标是 $(n_e+n_m)\bmod2$，这里写 $(n_m+n_e)\bmod2$，两者相同（加法可交换）， 但原文两处次序不同，抄写时容易抄错。$m=0$ 时两个奇点分别对应 $(\underline{2},\underline{1})$ 与 $(\underline{1},\underline{2})$， 荷 $(1,0)$ 与 $(1,1)$ 恰好落在两个不同的旋量中。
\par
\par\medskip\noindent\hangindent=3.2em\hangafter=1 \textbf{[198]}~\textbf{4.5}~见原文 PDF 第 75 页。这里 $\omega=e^{i\pi/6}$ 与式 (112) 后面的 $\omega=e^{2\pi i/[4(4-N_f)]}$ 不是同一个符号： $N_f=1$ 时后者给出 $e^{2\pi i/12}=e^{i\pi/6}$，两者一致。$\omega^2$ 是 ${\bf Z}_3$ 的本原根（$(\omega^2)^3=e^{2\pi i}=1$），故 $u\to\omega^4u=\omega^{-2}u$ 的 ${\bf Z}_3$ 作用把三个奇点循环置换， 电荷按 $(1,0)\to(1,1)\to(1,2)$ 轮换。此处符号与相位最易错，务必与 式 (112) 的 $\omega$ 定义统一。
\par
\par\medskip\noindent\hangindent=3.2em\hangafter=1 \textbf{[199]}~\textbf{4.6}~见原文 PDF 第 75--76 页。式 (116) 与第 3 节（原文 PDF 第 63 页）给出的 $y^2=(x^2-\Lambda^4)(x-u)$ 并不是同一个曲线族：本节把 $a$ 除以 2 而 $a_D$ 不变，故 $\tau=\partial a_D/\partial a$ 变为原来的 2 倍，曲线的模参数也随之改变（这正是第 65 页所说“电荷的重新标度改变 曲线形式”的含义）；两者由 $\tau\mapsto2\tau$ 相联系（相当于一个 2-等源 变换），并非简单的坐标重参数化。核对 $N_f=0$ 的周期与单值矩阵时，不要把 第 3 节的结果直接搬来用。
\par
\par\medskip\noindent\hangindent=3.2em\hangafter=1 \textbf{[200]}~\textbf{4.6}~见原文 PDF 第 76 页。“$p$ 与 $n$ 是模 2 的整数”实际是说 在 $\Gamma(2)$ 中 $p\equiv n\equiv0\pmod2$（故 $2n\equiv0\pmod4$、 $p/2$ 为整数）；于是换基后的单值群是 $\Gamma(2)$ 在此相似变换下的共轭， 即（按列矢量的写法）上右元为 4 的倍数的子群，原文记为 $\Gamma_0(4)$（按行矢量 $W$ 的写法即下左元为 4 的倍数）。可用旧矩阵 (104) 逐一代入核验：$M_1=\left(\begin{smallmatrix}1&0\\-2&1\end{smallmatrix}\right) \to\left(\begin{smallmatrix}1&0\\-1&1\end{smallmatrix}\right)$， $M_{-1}=\left(\begin{smallmatrix}-1&2\\-2&3\end{smallmatrix}\right) \to\left(\begin{smallmatrix}-1&4\\-1&3\end{smallmatrix}\right)$， $M_\infty=\left(\begin{smallmatrix}-1&2\\0&-1\end{smallmatrix}\right) \to\left(\begin{smallmatrix}-1&4\\0&-1\end{smallmatrix}\right)$， 且 $M_1M_{-1}=M_\infty$。原文把“matrices”误拼为“martices”，系排印笔误。
\par
\par\medskip\noindent\hangindent=3.2em\hangafter=1 \textbf{[201]}~\textbf{4.6}~见原文 PDF 第 76--77 页。此处是“$\lambda$ 平面绕行路径的 同伦类”论证：$\Gamma_1$ 连接 $\pm\lambda^{1/2}$，绕行一圈后两端同时 变号，恰好交换积分分支，故 $\omega_1$ 不变；$\Gamma_2$ 从 $\lambda^{1/2}$ 到 $1$，绕行一圈后其端点移到另一叶，路径同伦等价于 原来的 $\Gamma_2$ 加上 $\Gamma_1$，故 $\omega_2\to\omega_2+\omega_1$。 必须区分“以 $\lambda$ 为变量共轭于 $T$”与“以 $u$ 为变量共轭于 $T^n$”： 前者绕一圈、后者绕 $n$ 圈。图 6、图 7 的定向若画反，$\omega_2$ 的符号 会差一个负号，这是核对单值矩阵时最常见的错误来源。
\par
\par\medskip\noindent\hangindent=3.2em\hangafter=1 \textbf{[202]}~\textbf{4.6}~见原文 PDF 第 77 页。$u\to\infty$ 处的奇点“不稳定”指 多于两个分支点重合。用 $x=x'u,\ y=y'u^{3/2}$ 把分支点移到有限位置后， 判别式在 $(x',y')$ 变量下表现为 $\Delta\sim u^{-4}$，对应共轭于 $T^{-4}$ 的单值变换；但 $y=y'u^{3/2}$ 中的 $u^{3/2}$ 含有 $\sqrt u$， 绕 $u=0$ 一圈会给出额外的 $-1$，故回到原变量 $(x,y)$ 后单值变换是 $PT^{-4}$ 而不是 $T^{-4}$（$P=-1$ 即宇称）。漏掉这个负号、把结果写成 $T^{-4}$ 是最常见的错误。另外注意判别式的定义不含首项系数与整体符号， 与某些文献的约定相差常数因子。
\par
\par\medskip\noindent\hangindent=3.2em\hangafter=1 \textbf{[203]}~\textbf{4.6}~见原文 PDF 第 78 页。原文“massless magnetic monopoles with one unit of magnetic charge and $n_e$ units of magnetic charge”中第二个 “magnetic charge”应为“electric charge”：磁单极带一个单位的磁荷与 $n_e$ 个单位的电荷，这正是 $(n_m,n_e)=(1,n_e)$ 的含义。系原文笔误， 译文按正确含义（电荷）译出。另注意上表的单值矩阵之积给出 $M_\infty=PT^{N_f-4}$，可用 $S^2=-1$、$(ST)^3=1$ 逐项核验；它是后面 确定各曲线形式的主要约束。
\par
\par\medskip\noindent\hangindent=3.2em\hangafter=1 \textbf{[204]}~\textbf{4.6}~见原文 PDF 第 79 页。可用三次式 $x^3-ux^2+\tilde\Lambda^6$ 的判别式 $\Delta=4u^3\tilde\Lambda^6-27\tilde\Lambda^{12}$ 核验此式： $x^3+ax^2+bx+c$ 的判别式为 $a^2b^2-4b^3-4a^3c-27c^2+18abc$，代入 $a=-u,\ b=0,\ c=\tilde\Lambda^6$ 即得。注意原文的 $\Delta$ 定义不含 首项系数，与某些文献相差一个常数因子；零点位置（$u^3=\frac{27}{4} \tilde\Lambda^6$，三个零点相差 ${\bf Z}_3$ 相位）才是物理上重要的。
\par
\par\medskip\noindent\hangindent=3.2em\hangafter=1 \textbf{[205]}~\textbf{4.6}~见原文 PDF 第 79 页。注意 $(x^2-\tilde\Lambda_2^4)(x-u) = x^3-ux^2-\tilde\Lambda_2^4x+u\tilde\Lambda_2^4$ 确实具有上式的 “$(ax+bu)\Lambda_2^4$”形式：$a=-1,\ b=1$。所谓“判别式在每个有限奇点 处有二阶零点”即 $\Delta$ 有两个二重零点（对应两个奇点各两个无质量 单极），这正好确定两个系数 $a,b$。原文省略了这一代数步骤。
\par
\par\medskip\noindent\hangindent=3.2em\hangafter=1 \textbf{[206]}~\textbf{4.6}~见原文 PDF 第 79--80 页。原文只给出拟设与最终曲线，中间 “要求三次部分具有经典行为”和“经过一些重新标度”两步完全没有展开： 由 $F=x^3+(a\Lambda_3^2+cu)x^2+bu^2x$ 算判别式可得 $\Delta=b^2u^4[(a\Lambda_3^2+cu)^2-4bu^2]$，在 $u=0$ 处的四阶零点由 $b^2u^4$ 自动给出（$b\ne0$），这解释了下文“$b\not=0$”的条件；而 $x$、$y$ 的允许重参数化（含 $x$ 的平移/标度）会把拟设化为最终形式。 原文未展示这些步骤，若需复现请参看 [32] 的严格处理。
\par
\par\medskip\noindent\hangindent=3.2em\hangafter=1 \textbf{[207]}~\textbf{4.6}~见原文 PDF 第 80 页。这三条曲线在原书中全部带 nonumber 标注，因此不占公式编号，但它们与前面无质量曲线的关系可以 逐项核验：$m_i\to0$ 时 $N_f=1$ 给出 $-\frac1{64}\Lambda_1^6$（对应 $t=-1/64$），$N_f=2$ 给出 $(x^2-\frac1{64}\Lambda_2^4)(x-u)$（对应 $\tilde\Lambda_2^4=\frac1{64}\Lambda_2^4$），$N_f=3$ 给出 $-\frac1{64}\Lambda_3^2(x-u)^2$（对应 $\tilde\Lambda_3^2=-\frac1{64}\Lambda_3^2$，故 $\tilde\Lambda_3^2$ 的符号与 $\Lambda_3^2$ 相反）。核对符号时以这一极限为依据。
\par
\par\medskip\noindent\hangindent=3.2em\hangafter=1 \textbf{[208]}~\textbf{4.6}~见原文 PDF 第 80 页。$d\lambda/du$ 的系数 $\sqrt2/(8\pi)$ 中含因子 $\sqrt2$ 与 $8\pi$，与式 (109) 的 $1/2\pi$、式 (111) 的 $1/16\pi^2$ 相互关联；改动任何一处 归一化都会破坏周期 $a,a_D$ 与单值矩阵的一致性，抄写时必须逐字符核对。 另外原文“used by Seiberg and Witten to drive the curve”中的“drive” 系“derive”（导出）的拼写笔误；同页“varialbles”（应为 variables）、 “form the perturbative”（应为 from）等亦为排印笔误。全片类似的拼写 笔误还有 sigularities（应为 singularities，第 74--75 页）、 analoguous（应为 analogous，第 73 页）、representaion（应为 representation，第 73 页），译文均已按正确词义处理。留数正比于裸质量 是本节与前面 QED 讨论（$a_D$ 的平移）相呼应的关键结论。
\par
\par\medskip\noindent\hangindent=3.2em\hangafter=1 \textbf{[209]}~\textbf{致谢}~原文致谢中的 “Trieste Theory Group” 即 ICTP 的里雅斯特理论组；本讲义源自 1995 年 Trieste Summer School 的系列讲座（见原文 PDF 第 2 页引言第一段），故致谢对象与 讲座背景一致。
\par
